\documentclass[a4paper,12pt]{article}
\usepackage{float}
\usepackage[table]{xcolor} % Para colorir as linhas da tabela
%\usepackage{array} % Para melhorar a formatação de tabelas
\usepackage[brazil]{babel}
\usepackage[utf8]{inputenc}
\usepackage{diagbox}
\usepackage{amsmath}
\usepackage{amssymb}
\usepackage{multirow}
\usepackage{multicol}
\usepackage{graphicx}
%\usepackage{animate}
%\usepackage{tabularx}
\usepackage[font=footnotesize, labelfont=bf, listformat=empty]{caption}
\usepackage{geometry}
\usepackage{listings}
\usepackage{fontspec}
\usepackage{minted}
\usepackage{lastpage}
%\usepackage{subcaption}
\usepackage[skins,xparse,breakable]{tcolorbox}
%\usepackage{longtable}
\usepackage{titlesec}
\usepackage{fancyhdr}
\usepackage{ragged2e}
\usepackage{setspace}
\usepackage{booktabs}
\usepackage{enumitem}
\usepackage{minted}
\usepackage[colorlinks=true, allcolors=black]{hyperref}
\usepackage{tikz}

\usetikzlibrary{karnaugh}
\usetikzlibrary{positioning,arrows.meta,shapes.multipart,calc}

%\usetikzlibrary{arrows,shapes,automata,positioning,calc}
% Configurações de margens e espaçamento
\geometry{a4paper, left=2.5cm, right=2.5cm, top=2.5cm, bottom=2cm}
\onehalfspacing

% Configurações de títulos
\titleformat{\section}{\normalfont\Large\bfseries}{\thesection}{1em}{}
\titleformat{\subsection}{\normalfont\large\bfseries}{\thesubsection}{1em}{}
\titleformat{\subsubsection}
  {\normalfont\large}
  {\textbf{\thesubsubsection}}
  {0.25em}{}

\setcounter{tocdepth}{3}
\newcommand{\cod}[1]{\texttt{#1}}

\renewcommand{\listingscaption}{Código}
\newenvironment{code}{\captionsetup{type=listing}}{}

% Configurações de cores
\definecolor{LightSeaGreen}{rgb}{0.1255, 0.6980, 0.6667}
\definecolor{whitesmoke}{rgb}{0.96, 0.96, 0.96}
\definecolor{cinza}{RGB}{160, 160, 160}
\definecolor{cinza_claro}{RGB}{240, 240, 240} % Cinza claro
\definecolor{LogisimKMapColor0}{RGB}{128,0,0}
\definecolor{LogisimKMapColor1}{RGB}{230,25,75}
\definecolor{LogisimKMapColor2}{RGB}{250,190,190}
\definecolor{LogisimKMapColor3}{RGB}{170,110,40}
\definecolor{LogisimKMapColor4}{RGB}{245,130,48}
\definecolor{LogisimKMapColor5}{RGB}{255,215,180}
\definecolor{LogisimKMapColor6}{RGB}{128,128,0}
\definecolor{LogisimKMapColor7}{RGB}{255,255,25}
\definecolor{LogisimKMapColor8}{RGB}{210,245,60}
\definecolor{LogisimKMapColor9}{RGB}{0,0,128}
\definecolor{LogisimKMapColor10}{RGB}{145,30,180}
\definecolor{LogisimKMapColor11}{RGB}{60,180,175}
\definecolor{LogisimKMapColor12}{RGB}{0,130,203}
\definecolor{LogisimKMapColor13}{RGB}{230,190,255}
\definecolor{LogisimKMapColor14}{RGB}{170,255,195}
\definecolor{LogisimKMapColor15}{RGB}{240,50,230}

\newtcolorbox{codebox}[1]{
    title=#1,
    fonttitle=\small\bfseries,
    colback=blue!2.5!white,
    colframe=blue!62!white,
    arc=2pt,
    boxrule=0.8pt,
    left=4pt,
    right=4pt,
    top=4pt,
    bottom=4pt,
    breakable,
    enhanced jigsaw
}

\setminted[asm]{
    fontsize=\scriptsize,
    style=manni,
    linenos=true,
    breaklines=true,
    mathescape=true,
    bgcolor=cinza_claro,
    breakanywhere=true,
}

\newcommand{\hex}[1]{\textcolor{green!50!black}{#1}}

% Configurações de cabeçalho e rodapé
\pagestyle{fancy}
\fancyhf{}
\fancyhead[L]{\footnotesize{Relatório do Laboratório 3}}
\fancyhead[R]{\footnotesize{2025/2 - Turma 02}}
\fancyfoot[R]{\footnotesize{Página \hspace{0.05cm} \thepage \hspace{0.05cm} de \pageref{LastPage}}}
\fancyfoot[L]{\footnotesize{Organização e Arquitetura de Computadores - CIC0099}}

% Adiciona uma linha preta acima do rodapé
\renewcommand{\footrulewidth}{0.4pt} % Espessura da linha
\renewcommand{\footrule}{\vbox to 0pt{\hrule width \textwidth height \footrulewidth \vss}} % Desenha a linha

% Novo Comando para chamar a capa
\newcommand{\capa}{
    \begin{titlepage}
        \begin{multicols}{2}
            \begin{flushleft}
                
\includegraphics[width=0.45\linewidth]{Recursos/Imagens/UnB_logo.png}
            \end{flushleft}
            \columnbreak
            \begin{flushright}
                Universidade de Brasília \\
                Instituto de Exatas \\
                Departamento de Ciência da Computação
            \end{flushright}
        \end{multicols}
        \begin{center}
        \vspace{-20pt}
        \rule{\textwidth}{0.4pt}
        \end{center}
        \vspace{0.6cm}
        \begin{center}
            {\Huge \textbf{Relatório do Laboratório 3}} \\[0.5em]
            {\Huge \textbf{Processador Multiciclo}} \\ [0.5em]
            {\large CIC0099 - Organização e Arquitetura de Computadores} \\[0.25em]
            {\large Turma 02 - Prof. Dr. Ricardo Pezzuol Jacobi} \\
            \vfill
            {\Large \textbf{Grupo B4}} \\
            \begin{table}[H]
                \centering
                \begin{tabular}{l r}
                    Henrique Morcelles Salum & 232003008 \\
                    Athos Calixto Muniz & 211068261 \\
                    Gabriela Fernanda Rodrigues Costa & 180120859 \\
                    Cauê Araújo Euzebio & 211028195 \\
                    Lucas Silva Nóbrega & 180035096
                \end{tabular}
                \label{tab:placeholder}
            \end{table}
        \end{center}
    \end{titlepage}
}

\begin{document}

\capa

% Sumário
\newpage
\tableofcontents
\newpage

\section{Introdução}
O presente experimento tem como objetivo introduzir os alunos ao processo de projeto e implementação de sistemas digitais utilizando VHDL como Linguagem de Descrição de Hardware (HDL). A atividade busca promover o treinamento prático em modelagem, análise e síntese de circuitos digitais, utilizando o ambiente de desenvolvimento \textbf{Intel Quartus Prime v24.1} como ferramenta principal.

Como aplicação prática, propõe-se a implementação de um processador Multiciclo compatível com a ISA RV32I reduzida, contendo as instruções:\cod{add}, \cod{sub}, \cod{and}, \cod{or}, \cod{slt}, \cod{lw}, \cod{sw}, \cod{beq}, \cod{jal}, \cod{jalr}, \cod{addi} e \cod{lui}.

A implementação será validada por meio de simulação funcional e temporal, com análise de \textit{timing} e frequência máxima de operação, permitindo compreender, de forma integrada, o funcionamento de uma arquitetura RISC-V.


\section{Métodos e Análise}

\subsection{Experimento}
\begin{tcolorbox}[title=Enunciado, colback=blue!5!white, colframe=blue!75!black]
Implemente o processador Multiciclo com ISA Reduzida com as instruções: \cod{add}, \cod{sub}, \cod{and}, \cod{or}, \cod{slt}, \cod{lw},
\cod{sw}, \cod{beq}, \cod{jal}, e ainda as instruções \cod{jalr}, \cod{addi} e \cod{lui}.
\end{tcolorbox}
Para realizar o solicitado, é-nos fornecido o esquemático de uma implementação do Multiciclo com ISA reduzida. Nesse esquemático, porém, foi necessário realizar alterações no caminho de dados, que serão devidamente justificadas adiante.

\begin{figure}[H]
    \centering
    % Important: If latex complains about unicode characters,
% please use "\usepackage[utf8x]{inputenc}" in your preamble
% You can change the size of the picture by putting it into the construct:
% 1) \resizebox{10cm}{!}{"below picture"} to scale horizontally to 10 cm
% 2) \resizebox{!}{15cm}{"below picture"} to scale vertically to 15 cm
% 3) \resizebox{10cm}{15cm}{"below picture"} a combination of above two
% It is not recomended to use the scale option of the tikzpicture environment.
\resizebox{\textwidth}{!}{
\begin{tikzpicture}[x=1pt,y=-1pt,line cap=rect]
\def\logisimfontA#1{{\fontspec{CMU Serif}#1}} % cmr equivalent
\def\logisimfontB#1{\fontspec{CMU Bright}#1}
\def\logisimfontC#1{\fontspec{Arial Rounded MT Bold}#1}
\definecolor{custcol_87_87_87}{RGB}{135, 135, 135}
\definecolor{custcol_0_0_0}{RGB}{0, 0, 0}
\definecolor{custcol_ff_ff_ff}{RGB}{255, 255, 255}
\draw [line width=4.0pt, custcol_0_0_0 ]  (735.0,175.0) -- (785.0,175.0) ;
\draw [line width=3.0pt, custcol_0_0_0 ]  (165.0,165.0) -- (165.0,175.0) ;
\draw [line width=3.0pt, custcol_0_0_0 ]  (45.0,45.0) -- (45.0,55.0) ;
\draw [line width=3.0pt, custcol_0_0_0 ]  (295.0,95.0) -- (295.0,105.0) ;
\draw [line width=3.0pt, custcol_0_0_0 ]  (385.0,145.0) -- (385.0,155.0) ;
\draw [line width=4.0pt, custcol_0_0_0 ]  (805.0,365.0) -- (805.0,385.0) ;
\draw [line width=4.0pt, custcol_0_0_0 ]  (185.0,135.0) -- (205.0,135.0) ;
\draw [line width=4.0pt, custcol_0_0_0 ]  (825.0,325.0) -- (855.0,325.0) ;
\draw [line width=3.0pt, custcol_0_0_0 ]  (915.0,75.0) -- (925.0,75.0) ;
\draw [line width=3.0pt, custcol_0_0_0 ]  (625.0,35.0) -- (625.0,45.0) ;
\draw [line width=3.0pt, custcol_0_0_0 ]  (1015.0,185.0) -- (1015.0,195.0) ;
\draw [line width=4.0pt, custcol_0_0_0 ]  (735.0,295.0) -- (755.0,295.0) -- (785.0,295.0) ;
\draw [line width=4.0pt, custcol_0_0_0 ]  (205.0,295.0) -- (165.0,295.0) -- (165.0,525.0) -- (755.0,525.0) -- (755.0,295.0) ;
\draw [line width=4.0pt, custcol_0_0_0 ]  (485.0,135.0) -- (495.0,135.0) ;
\draw [line width=3.0pt, custcol_0_0_0 ]  (1015.0,65.0) -- (1025.0,65.0) ;
\draw [line width=4.0pt, custcol_0_0_0 ]  (475.0,425.0) -- (465.0,425.0) -- (465.0,505.0) -- (725.0,505.0) -- (725.0,405.0) -- (675.0,405.0) ;
\draw [line width=4.0pt, custcol_0_0_0 ]  (915.0,285.0) -- (915.0,335.0) ;
\draw [line width=3.0pt, custcol_0_0_0 ]  (645.0,145.0) -- (645.0,155.0) ;
\draw [line width=3.0pt, custcol_0_0_0 ]  (235.0,95.0) -- (235.0,105.0) ;
\draw [line width=4.0pt, custcol_0_0_0 ]  (495.0,435.0) -- (495.0,445.0) ;
\draw [line width=4.0pt, custcol_0_0_0 ]  (725.0,405.0) -- (725.0,335.0) -- (785.0,335.0) ;
\draw [line width=4.0pt, custcol_0_0_0 ]  (805.0,225.0) -- (805.0,245.0) ;
\draw [line width=4.0pt, custcol_0_0_0 ]  (765.0,225.0) -- (765.0,215.0) -- (785.0,215.0) ;
\draw [line width=4.0pt, custcol_0_0_0 ]  (325.0,215.0) -- (345.0,215.0) ;
\draw [line width=4.0pt, custcol_0_0_0 ]  (485.0,295.0) -- (505.0,295.0) ;
\draw [line width=4.0pt, custcol_0_0_0 ]  (145.0,115.0) -- (115.0,115.0) -- (115.0,375.0) -- (455.0,375.0) -- (455.0,385.0) -- (475.0,385.0) ;
\draw [line width=4.0pt, custcol_0_0_0 ]  (825.0,185.0) -- (855.0,185.0) ;
\draw [line width=3.0pt, custcol_0_0_0 ]  (915.0,225.0) -- (915.0,95.0) -- (925.0,95.0) ;
\draw [line width=4.0pt, custcol_0_0_0 ]  (485.0,175.0) -- (555.0,175.0) ;
\draw [line width=4.0pt, custcol_0_0_0 ]  (485.0,215.0) -- (555.0,215.0) ;
\draw [line width=4.0pt, custcol_0_0_0 ]  (485.0,255.0) -- (555.0,255.0) ;
\draw [line width=4.0pt, custcol_0_0_0 ]  (665.0,65.0) -- (775.0,65.0) -- (775.0,155.0) -- (785.0,155.0) ;
\draw [line width=3.0pt, custcol_0_0_0 ]  (775.0,315.0) -- (785.0,315.0) ;
\draw [line width=4.0pt, custcol_0_0_0 ]  (1025.0,255.0) -- (1075.0,255.0) -- (1075.0,545.0) -- (135.0,545.0) -- (135.0,365.0) -- (135.0,155.0) -- (145.0,155.0) ;
\draw [line width=4.0pt, custcol_0_0_0 ]  (995.0,175.0) -- (975.0,175.0) -- (975.0,225.0) -- (1075.0,225.0) -- (1075.0,255.0) ;
\draw [line width=4.0pt, custcol_0_0_0 ]  (1035.0,155.0) -- (1075.0,155.0) -- (1075.0,5.0) -- (5.0,5.0) -- (5.0,115.0) -- (25.0,115.0) ;
\draw [line width=4.0pt, custcol_0_0_0 ]  (465.0,215.0) -- (455.0,215.0) -- (455.0,325.0) -- (555.0,325.0) -- (555.0,405.0) -- (575.0,405.0) ;
\draw [line width=4.0pt, custcol_0_0_0 ]  (785.0,355.0) -- (765.0,355.0) -- (765.0,375.0) ;
\draw [line width=4.0pt, custcol_0_0_0 ]  (555.0,65.0) -- (555.0,105.0) -- (765.0,105.0) -- (765.0,195.0) -- (785.0,195.0) ;
\draw [line width=4.0pt, custcol_0_0_0 ]  (935.0,255.0) -- (955.0,255.0) -- (975.0,255.0) ;
\draw [line width=4.0pt, custcol_0_0_0 ]  (675.0,175.0) -- (695.0,175.0) ;
\draw [line width=4.0pt, custcol_0_0_0 ]  (675.0,295.0) -- (695.0,295.0) ;
\draw [line width=4.0pt, custcol_0_0_0 ]  (365.0,215.0) -- (345.0,215.0) -- (345.0,445.0) -- (365.0,445.0) ;
\draw [line width=4.0pt, custcol_0_0_0 ]  (435.0,215.0) -- (455.0,215.0) ;
\draw [line width=4.0pt, custcol_0_0_0 ]  (135.0,365.0) -- (475.0,365.0) ;
\draw [line width=4.0pt, custcol_0_0_0 ]  (475.0,405.0) -- (455.0,405.0) -- (455.0,445.0) -- (435.0,445.0) ;
\draw [line width=4.0pt, custcol_0_0_0 ]  (615.0,65.0) -- (555.0,65.0) -- (115.0,65.0) -- (115.0,115.0) -- (95.0,115.0) ;
\draw [line width=4.0pt, custcol_0_0_0 ]  (995.0,135.0) -- (955.0,135.0) -- (955.0,255.0) ;
\draw [line width=4.0pt, custcol_0_0_0 ]  (515.0,395.0) -- (545.0,395.0) -- (545.0,295.0) -- (555.0,295.0) ;
\fill [line width=4.0pt, custcol_0_0_0]  (955.0,255.0) ellipse (6.0 and 6.0 );
\fill [line width=4.0pt, custcol_0_0_0]  (135.0,365.0) ellipse (6.0 and 6.0 );
\fill [line width=4.0pt, custcol_0_0_0]  (455.0,215.0) ellipse (6.0 and 6.0 );
\fill [line width=4.0pt, custcol_0_0_0]  (1075.0,255.0) ellipse (6.0 and 6.0 );
\fill [line width=4.0pt, custcol_0_0_0]  (115.0,115.0) ellipse (6.0 and 6.0 );
\fill [line width=4.0pt, custcol_0_0_0]  (725.0,405.0) ellipse (6.0 and 6.0 );
\fill [line width=4.0pt, custcol_0_0_0]  (555.0,65.0) ellipse (6.0 and 6.0 );
\fill [line width=4.0pt, custcol_0_0_0]  (755.0,295.0) ellipse (6.0 and 6.0 );
\fill [line width=4.0pt, custcol_0_0_0]  (345.0,215.0) ellipse (6.0 and 6.0 );
\draw [line width=4.0pt, custcol_0_0_0 ]  (145.0,95.0) -- (145.0,175.0) -- (185.0,155.0) -- (185.0,115.0) -- cycle;
\logisimfontB{\fontsize{12pt}{12pt}\selectfont\node[inner sep=0, outer sep=0, custcol_87_87_87, anchor=base west] at  (152.0,119.0)  {0};}
\logisimfontC{\fontsize{14pt}{14pt}\fontseries{bx}\selectfont\node[inner sep=0, outer sep=0, custcol_0_0_0, anchor=base west] at  (151.0,135.0)  {Mux};}
\logisimfontC{\fontsize{10pt}{10pt}\fontseries{bx}\selectfont\node[inner sep=0, outer sep=0, custcol_0_0_0, anchor=base west] at  (156.0,144.0)  {2x1};}
\fill [line width=1.0pt, custcol_0_0_0]  (145.0,115.0) ellipse (2.0 and 2.0 );
\fill [line width=1.0pt, custcol_0_0_0]  (145.0,155.0) ellipse (2.0 and 2.0 );
\fill [line width=1.0pt, custcol_0_0_0]  (165.0,165.0) ellipse (2.0 and 2.0 );
\fill [line width=1.0pt, custcol_0_0_0]  (185.0,135.0) ellipse (2.0 and 2.0 );
\logisimfontC{\fontsize{20pt}{20pt}\fontseries{bx}\selectfont\node[inner sep=0, outer sep=0, custcol_0_0_0, anchor=base west] at  (570.0,235.0)  {Register};}
\logisimfontC{\fontsize{20pt}{20pt}\fontseries{bx}\selectfont\node[inner sep=0, outer sep=0, custcol_0_0_0, anchor=base west] at  (588.0,254.0)  {Bank};}
\draw [line width=1.0pt, custcol_0_0_0 ]  (640.0,315.0) -- (645.0,301.0) -- (650.0,315.0) ;
\begin{pgfpicture}
   \begin{pgfmagnify}{1pt}{-1pt}
      \pgfsetrectcap
      \pgfsetcornersarced{\pgfpoint{10.0}{10.0}}
      \pgfsetlinewidth{4.0}
      \color{custcol_0_0_0}
      \pgfsetfillopacity{1.0}
      \pgfpathrectanglecorners{\pgfpoint{555.0}{155.0}}{\pgfpoint{675.0}{315.0}}
      \pgfusepath{stroke}
   \end{pgfmagnify}
\end{pgfpicture}
\fill [line width=1.0pt, custcol_0_0_0]  (555.0,175.0) ellipse (2.0 and 2.0 );
\fill [line width=1.0pt, custcol_0_0_0]  (555.0,215.0) ellipse (2.0 and 2.0 );
\fill [line width=1.0pt, custcol_0_0_0]  (555.0,255.0) ellipse (2.0 and 2.0 );
\fill [line width=1.0pt, custcol_0_0_0]  (555.0,295.0) ellipse (2.0 and 2.0 );
\fill [line width=1.0pt, custcol_0_0_0]  (645.0,155.0) ellipse (2.0 and 2.0 );
\fill [line width=1.0pt, custcol_0_0_0]  (645.0,315.0) ellipse (2.0 and 2.0 );
\fill [line width=1.0pt, custcol_0_0_0]  (675.0,175.0) ellipse (2.0 and 2.0 );
\fill [line width=1.0pt, custcol_0_0_0]  (675.0,295.0) ellipse (2.0 and 2.0 );
\logisimfontC{\fontsize{8pt}{8pt}\fontseries{bx}\selectfont\node[inner sep=0, outer sep=0, custcol_0_0_0, anchor=base west] at  (275.0,88.0)  {memRead};}
\draw [line width=2.0pt, custcol_0_0_0 ]  (271.0,77.0) -- (319.0,77.0) -- (319.0,92.0) -- (305.0,92.0) -- (295.0,96.0) -- (285.0,92.0) -- (271.0,92.0) -- cycle;
\fill [line width=2.0pt, custcol_0_0_0]  (295.0,95.0) ellipse (2.0 and 2.0 );
\logisimfontC{\fontsize{20pt}{20pt}\fontseries{bx}\selectfont\node[inner sep=0, outer sep=0, custcol_0_0_0, anchor=base west] at  (223.0,221.0)  {Memory};}
\draw [line width=4.0pt, custcol_0_0_0 ]  (205.0,105.0) -- (324.0,105.0) ;
\draw [line width=4.0pt, custcol_0_0_0 ]  (325.0,105.0) -- (325.0,324.0) ;
\draw [line width=4.0pt, custcol_0_0_0 ]  (325.0,325.0) -- (206.0,325.0) ;
\draw [line width=4.0pt, custcol_0_0_0 ]  (205.0,325.0) -- (205.0,106.0) ;
\draw [line width=1.0pt, custcol_0_0_0 ]  (289.0,324.0) -- (295.0,311.0) -- (301.0,324.0) ;
\logisimfontB{\fontsize{14pt}{14pt}\selectfont\node[inner sep=0, outer sep=0, custcol_87_87_87, anchor=base west] at  (214.0,139.0)  {address};}
\logisimfontB{\fontsize{14pt}{14pt}\selectfont\node[inner sep=0, outer sep=0, custcol_87_87_87, anchor=base west] at  (215.0,299.0)  {data};}
\draw [line width=2.0pt, custcol_0_0_0 ]  (203.0,105.0) -- (200.0,110.0) -- (201.0,329.0) -- (329.0,329.0) -- (330.0,110.0) -- (327.0,105.0) ;
\draw [line width=1.0pt, custcol_0_0_0 ]  (203.0,104.0) -- (197.0,110.0) -- (199.0,331.0) -- (331.0,331.0) -- (333.0,110.0) -- (327.0,104.0) ;
\fill [line width=1.0pt, custcol_0_0_0]  (205.0,135.0) ellipse (2.0 and 2.0 );
\fill [line width=1.0pt, custcol_0_0_0]  (205.0,295.0) ellipse (2.0 and 2.0 );
\fill [line width=1.0pt, custcol_0_0_0]  (235.0,105.0) ellipse (2.0 and 2.0 );
\fill [line width=1.0pt, custcol_0_0_0]  (295.0,105.0) ellipse (2.0 and 2.0 );
\fill [line width=1.0pt, custcol_0_0_0]  (295.0,325.0) ellipse (2.0 and 2.0 );
\fill [line width=1.0pt, custcol_0_0_0]  (325.0,215.0) ellipse (2.0 and 2.0 );
\logisimfontC{\fontsize{7pt}{7pt}\fontseries{bx}\selectfont\node[inner sep=0, outer sep=0, custcol_0_0_0, anchor=base west] at  (510.0,298.0)  {FUNCT};}
\draw [line width=2.0pt, custcol_0_0_0 ]  (509.0,287.0) -- (537.0,287.0) -- (537.0,303.0) -- (509.0,303.0) -- (505.0,295.0) -- cycle;
\fill [line width=2.0pt, custcol_0_0_0]  (505.0,295.0) ellipse (2.0 and 2.0 );
\logisimfontC{\fontsize{10pt}{10pt}\fontseries{bx}\selectfont\node[inner sep=0, outer sep=0, custcol_0_0_0, anchor=base west] at  (623.0,137.0)  {regWrite};}
\draw [line width=2.0pt, custcol_0_0_0 ]  (620.0,124.0) -- (671.0,124.0) -- (671.0,142.0) -- (655.0,142.0) -- (645.0,146.0) -- (635.0,142.0) -- (620.0,142.0) -- cycle;
\fill [line width=2.0pt, custcol_0_0_0]  (645.0,145.0) ellipse (2.0 and 2.0 );
\logisimfontC{\fontsize{10pt}{10pt}\fontseries{bx}\selectfont\node[inner sep=0, outer sep=0, custcol_0_0_0, anchor=base west] at  (1000.0,210.0)  {PCsrc};}
\draw [line width=2.0pt, custcol_0_0_0 ]  (997.0,199.0) -- (1005.0,199.0) -- (1015.0,195.0) -- (1025.0,199.0) -- (1034.0,199.0) -- (1034.0,216.0) -- (997.0,216.0) -- cycle;
\fill [line width=2.0pt, custcol_0_0_0]  (1015.0,195.0) ellipse (2.0 and 2.0 );
\draw [line width=2.0pt, custcol_0_0_0] (940.0,100.0) arc (90.0:-90.0:15.0 and 15.0 );
\draw [line width=2.0pt, custcol_0_0_0 ]  (940.0,70.0) -- (926.0,70.0) -- (926.0,100.0) -- (940.0,100.0) ;
\draw [line width=2.0pt, custcol_0_0_0 ]  (975.0,235.0) -- (1024.0,235.0) ;
\draw [line width=2.0pt, custcol_0_0_0 ]  (1025.0,235.0) -- (1025.0,274.0) ;
\draw [line width=2.0pt, custcol_0_0_0 ]  (1025.0,275.0) -- (976.0,275.0) ;
\draw [line width=2.0pt, custcol_0_0_0 ]  (975.0,275.0) -- (975.0,236.0) ;
\draw [line width=1.0pt, custcol_0_0_0 ]  (1026.0,239.0) -- (1029.0,239.0) -- (1029.0,279.0) -- (980.0,279.0) -- (980.0,276.0) ;
\draw [line width=1.0pt, custcol_0_0_0 ]  (1013.0,275.0) -- (1015.0,268.0) -- (1017.0,275.0) ;
\draw [line width=1.0pt, custcol_0_0_0 ]  (1030.0,242.0) -- (1032.0,242.0) -- (1032.0,282.0) -- (985.0,282.0) -- (985.0,280.0) ;
\logisimfontC{\fontsize{11pt}{11pt}\fontseries{bx}\selectfont\node[inner sep=0, outer sep=0, custcol_0_0_0, anchor=base west] at  (978.0,259.0)  {ALUout};}
\fill [line width=1.0pt, custcol_0_0_0]  (975.0,255.0) ellipse (2.0 and 2.0 );
\fill [line width=1.0pt, custcol_0_0_0]  (1015.0,275.0) ellipse (2.0 and 2.0 );
\fill [line width=1.0pt, custcol_0_0_0]  (1025.0,255.0) ellipse (2.0 and 2.0 );
\draw [line width=2.0pt, custcol_0_0_0 ]  (435.0,160.0) -- (442.0,160.0) -- (442.0,282.0) -- (370.0,282.0) -- (370.0,275.0) ;
\draw [line width=1.0pt, custcol_0_0_0 ]  (409.0,275.0) -- (415.0,265.0) -- (421.0,275.0) ;
\draw [line width=4.0pt, custcol_0_0_0 ]  (365.0,155.0) -- (434.0,155.0) ;
\draw [line width=4.0pt, custcol_0_0_0 ]  (435.0,155.0) -- (435.0,274.0) ;
\draw [line width=4.0pt, custcol_0_0_0 ]  (435.0,275.0) -- (366.0,275.0) ;
\draw [line width=4.0pt, custcol_0_0_0 ]  (365.0,275.0) -- (365.0,156.0) ;
\draw [line width=2.0pt, custcol_0_0_0 ]  (443.0,165.0) -- (447.0,165.0) -- (447.0,287.0) -- (375.0,287.0) -- (375.0,283.0) ;
\logisimfontC{\fontsize{11pt}{11pt}\fontseries{bx}\selectfont\node[inner sep=0, outer sep=0, custcol_0_0_0, anchor=base west] at  (368.0,214.0)  {Instruction};}
\logisimfontC{\fontsize{11pt}{11pt}\fontseries{bx}\selectfont\node[inner sep=0, outer sep=0, custcol_0_0_0, anchor=base west] at  (375.0,224.0)  {Register};}
\fill [line width=1.0pt, custcol_0_0_0]  (365.0,215.0) ellipse (2.0 and 2.0 );
\fill [line width=1.0pt, custcol_0_0_0]  (385.0,155.0) ellipse (2.0 and 2.0 );
\fill [line width=1.0pt, custcol_0_0_0]  (415.0,275.0) ellipse (2.0 and 2.0 );
\fill [line width=1.0pt, custcol_0_0_0]  (435.0,215.0) ellipse (2.0 and 2.0 );
\draw [line width=2.0pt, custcol_0_0_0 ]  (695.0,155.0) -- (734.0,155.0) ;
\draw [line width=2.0pt, custcol_0_0_0 ]  (735.0,155.0) -- (735.0,194.0) ;
\draw [line width=2.0pt, custcol_0_0_0 ]  (735.0,195.0) -- (696.0,195.0) ;
\draw [line width=2.0pt, custcol_0_0_0 ]  (695.0,195.0) -- (695.0,156.0) ;
\draw [line width=1.0pt, custcol_0_0_0 ]  (723.0,195.0) -- (725.0,188.0) -- (727.0,195.0) ;
\draw [line width=1.0pt, custcol_0_0_0 ]  (736.0,159.0) -- (739.0,159.0) -- (739.0,199.0) -- (699.0,199.0) -- (699.0,196.0) ;
\draw [line width=1.0pt, custcol_0_0_0 ]  (739.0,162.0) -- (742.0,162.0) -- (742.0,202.0) -- (702.0,202.0) -- (702.0,200.0) ;
\logisimfontC{\fontsize{22pt}{22pt}\fontseries{bx}\selectfont\node[inner sep=0, outer sep=0, custcol_0_0_0, anchor=base west] at  (707.0,184.0)  {A};}
\fill [line width=1.0pt, custcol_0_0_0]  (695.0,175.0) ellipse (2.0 and 2.0 );
\fill [line width=1.0pt, custcol_0_0_0]  (725.0,195.0) ellipse (2.0 and 2.0 );
\fill [line width=1.0pt, custcol_0_0_0]  (735.0,175.0) ellipse (2.0 and 2.0 );
\draw [line width=4.0pt, custcol_0_0_0 ]  (995.0,115.0) -- (995.0,195.0) -- (1035.0,175.0) -- (1035.0,135.0) -- cycle;
\logisimfontB{\fontsize{12pt}{12pt}\selectfont\node[inner sep=0, outer sep=0, custcol_87_87_87, anchor=base west] at  (1002.0,139.0)  {0};}
\logisimfontC{\fontsize{14pt}{14pt}\fontseries{bx}\selectfont\node[inner sep=0, outer sep=0, custcol_0_0_0, anchor=base west] at  (1001.0,155.0)  {Mux};}
\logisimfontC{\fontsize{10pt}{10pt}\fontseries{bx}\selectfont\node[inner sep=0, outer sep=0, custcol_0_0_0, anchor=base west] at  (1006.0,164.0)  {2x1};}
\fill [line width=1.0pt, custcol_0_0_0]  (995.0,135.0) ellipse (2.0 and 2.0 );
\fill [line width=1.0pt, custcol_0_0_0]  (995.0,175.0) ellipse (2.0 and 2.0 );
\fill [line width=1.0pt, custcol_0_0_0]  (1015.0,185.0) ellipse (2.0 and 2.0 );
\fill [line width=1.0pt, custcol_0_0_0]  (1035.0,155.0) ellipse (2.0 and 2.0 );
\logisimfontC{\fontsize{10pt}{10pt}\fontseries{bx}\selectfont\node[inner sep=0, outer sep=0, custcol_0_0_0, anchor=base west] at  (895.0,350.0)  {ALUCtrl};}
\draw [line width=2.0pt, custcol_0_0_0 ]  (892.0,339.0) -- (905.0,339.0) -- (915.0,335.0) -- (925.0,339.0) -- (939.0,339.0) -- (939.0,356.0) -- (892.0,356.0) -- cycle;
\fill [line width=2.0pt, custcol_0_0_0]  (915.0,335.0) ellipse (2.0 and 2.0 );
\logisimfontC{\fontsize{8pt}{8pt}\fontseries{bx}\selectfont\node[inner sep=0, outer sep=0, custcol_0_0_0, anchor=base west] at  (501.0,138.0)  {OPCODE};}
\draw [line width=2.0pt, custcol_0_0_0 ]  (499.0,127.0) -- (539.0,127.0) -- (539.0,143.0) -- (499.0,143.0) -- (495.0,135.0) -- cycle;
\fill [line width=2.0pt, custcol_0_0_0]  (495.0,135.0) ellipse (2.0 and 2.0 );
\logisimfontC{\fontsize{18pt}{18pt}\fontseries{bx}\selectfont\node[inner sep=0, outer sep=0, custcol_0_0_0, anchor=base west] at  (890.0,262.0)  {ALU};}
\logisimfontB{\fontsize{10pt}{10pt}\selectfont\node[inner sep=0, outer sep=0, custcol_87_87_87, anchor=base west] at  (900.0,238.0)  {zero};}
\draw [line width=4.0pt, custcol_0_0_0 ]  (855.0,165.0) -- (855.0,205.0) -- (885.0,234.0) -- (885.0,275.0) -- (855.0,305.0) -- (855.0,345.0) -- (935.0,265.0) -- (935.0,245.0) -- cycle;
\fill [line width=1.0pt, custcol_0_0_0]  (855.0,185.0) ellipse (2.0 and 2.0 );
\fill [line width=1.0pt, custcol_0_0_0]  (855.0,325.0) ellipse (2.0 and 2.0 );
\fill [line width=1.0pt, custcol_0_0_0]  (915.0,225.0) ellipse (2.0 and 2.0 );
\fill [line width=1.0pt, custcol_0_0_0]  (915.0,285.0) ellipse (2.0 and 2.0 );
\fill [line width=1.0pt, custcol_0_0_0]  (935.0,255.0) ellipse (2.0 and 2.0 );
\draw [line width=3.0pt, custcol_0_0_0 ]  (765.0,225.0) -- (765.0,230.0) ;
\draw [line width=1.0pt, custcol_0_0_0 ]  (773.0,231.0) -- (757.0,231.0) ;
\draw [line width=1.0pt, custcol_0_0_0 ]  (770.0,234.0) -- (760.0,234.0) ;
\draw [line width=1.0pt, custcol_0_0_0 ]  (767.0,237.0) -- (763.0,237.0) ;
\fill [line width=1.0pt, custcol_0_0_0]  (765.0,225.0) ellipse (2.0 and 2.0 );
\logisimfontC{\fontsize{10pt}{10pt}\fontseries{bx}\selectfont\node[inner sep=0, outer sep=0, custcol_0_0_0, anchor=base west] at  (783.0,260.0)  {ALUsrcA};}
\draw [line width=2.0pt, custcol_0_0_0 ]  (780.0,249.0) -- (795.0,249.0) -- (805.0,245.0) -- (815.0,249.0) -- (831.0,249.0) -- (831.0,266.0) -- (780.0,266.0) -- cycle;
\fill [line width=2.0pt, custcol_0_0_0]  (805.0,245.0) ellipse (2.0 and 2.0 );
\logisimfontB{\fontsize{14pt}{14pt}\selectfont\node[inner sep=0, outer sep=0, custcol_87_87_87, anchor=base west] at  (481.0,370.0)  {0};}
\draw [line width=4.0pt, custcol_0_0_0 ]  (475.0,345.0) -- (475.0,445.0) -- (515.0,425.0) -- (515.0,365.0) -- cycle;
\logisimfontC{\fontsize{14pt}{14pt}\fontseries{bx}\selectfont\node[inner sep=0, outer sep=0, custcol_0_0_0, anchor=base west] at  (481.0,395.0)  {Mux};}
\logisimfontC{\fontsize{14pt}{14pt}\fontseries{bx}\selectfont\node[inner sep=0, outer sep=0, custcol_0_0_0, anchor=base west] at  (484.0,406.0)  {4x1};}
\fill [line width=1.0pt, custcol_0_0_0]  (475.0,365.0) ellipse (2.0 and 2.0 );
\fill [line width=1.0pt, custcol_0_0_0]  (475.0,385.0) ellipse (2.0 and 2.0 );
\fill [line width=1.0pt, custcol_0_0_0]  (475.0,405.0) ellipse (2.0 and 2.0 );
\fill [line width=1.0pt, custcol_0_0_0]  (475.0,425.0) ellipse (2.0 and 2.0 );
\fill [line width=1.0pt, custcol_0_0_0]  (495.0,435.0) ellipse (2.0 and 2.0 );
\fill [line width=1.0pt, custcol_0_0_0]  (515.0,395.0) ellipse (2.0 and 2.0 );
\draw [line width=3.0pt, custcol_0_0_0 ]  (485.0,135.0) -- (471.0,135.0) ;
\draw [line width=3.0pt, custcol_0_0_0 ]  (485.0,175.0) -- (471.0,175.0) ;
\draw [line width=3.0pt, custcol_0_0_0 ]  (485.0,215.0) -- (471.0,215.0) ;
\draw [line width=3.0pt, custcol_0_0_0 ]  (485.0,255.0) -- (471.0,255.0) ;
\draw [line width=3.0pt, custcol_0_0_0 ]  (485.0,295.0) -- (471.0,295.0) ;
\draw [line width=5.0pt, custcol_0_0_0 ]  (466.0,215.0) -- (471.0,215.0) ;
\draw [line width=5.0pt, custcol_0_0_0 ]  (471.0,136.0) -- (471.0,294.0) ;
\logisimfontA{\fontsize{7pt}{7pt}\selectfont\node[inner sep=0, outer sep=0, custcol_0_0_0, anchor=base west] at  (474.0,132.0)  {6-0};}
\logisimfontA{\fontsize{7pt}{7pt}\selectfont\node[inner sep=0, outer sep=0, custcol_0_0_0, anchor=base west] at  (474.0,172.0)  {19-15};}
\logisimfontA{\fontsize{7pt}{7pt}\selectfont\node[inner sep=0, outer sep=0, custcol_0_0_0, anchor=base west] at  (474.0,212.0)  {24-20};}
\logisimfontA{\fontsize{7pt}{7pt}\selectfont\node[inner sep=0, outer sep=0, custcol_0_0_0, anchor=base west] at  (474.0,252.0)  {11-7};}
\logisimfontA{\fontsize{7pt}{7pt}\selectfont\node[inner sep=0, outer sep=0, custcol_0_0_0, anchor=base west] at  (474.0,292.0)  {30,14-12};}
\fill [line width=5.0pt, custcol_0_0_0]  (465.0,215.0) ellipse (2.0 and 2.0 );
\fill [line width=5.0pt, custcol_0_0_0]  (485.0,135.0) ellipse (2.0 and 2.0 );
\fill [line width=5.0pt, custcol_0_0_0]  (485.0,175.0) ellipse (2.0 and 2.0 );
\fill [line width=5.0pt, custcol_0_0_0]  (485.0,215.0) ellipse (2.0 and 2.0 );
\fill [line width=5.0pt, custcol_0_0_0]  (485.0,255.0) ellipse (2.0 and 2.0 );
\fill [line width=5.0pt, custcol_0_0_0]  (485.0,295.0) ellipse (2.0 and 2.0 );
\logisimfontC{\fontsize{10pt}{10pt}\fontseries{bx}\selectfont\node[inner sep=0, outer sep=0, custcol_0_0_0, anchor=base west] at  (1030.0,69.0)  {PCctrl};}
\draw [line width=2.0pt, custcol_0_0_0 ]  (1029.0,56.0) -- (1066.0,56.0) -- (1066.0,75.0) -- (1029.0,75.0) -- (1025.0,65.0) -- cycle;
\fill [line width=2.0pt, custcol_0_0_0]  (1025.0,65.0) ellipse (2.0 and 2.0 );
\logisimfontC{\fontsize{10pt}{10pt}\fontseries{bx}\selectfont\node[inner sep=0, outer sep=0, custcol_0_0_0, anchor=base west] at  (155.0,190.0)  {IorD};}
\draw [line width=2.0pt, custcol_0_0_0 ]  (152.0,179.0) -- (155.0,179.0) -- (165.0,175.0) -- (175.0,179.0) -- (179.0,179.0) -- (179.0,196.0) -- (152.0,196.0) -- cycle;
\fill [line width=2.0pt, custcol_0_0_0]  (165.0,175.0) ellipse (2.0 and 2.0 );
\logisimfontB{\fontsize{14pt}{14pt}\selectfont\node[inner sep=0, outer sep=0, custcol_87_87_87, anchor=base west] at  (791.0,300.0)  {0};}
\draw [line width=4.0pt, custcol_0_0_0 ]  (785.0,275.0) -- (785.0,375.0) -- (825.0,355.0) -- (825.0,295.0) -- cycle;
\logisimfontC{\fontsize{14pt}{14pt}\fontseries{bx}\selectfont\node[inner sep=0, outer sep=0, custcol_0_0_0, anchor=base west] at  (791.0,325.0)  {Mux};}
\logisimfontC{\fontsize{14pt}{14pt}\fontseries{bx}\selectfont\node[inner sep=0, outer sep=0, custcol_0_0_0, anchor=base west] at  (794.0,336.0)  {4x1};}
\fill [line width=1.0pt, custcol_0_0_0]  (785.0,295.0) ellipse (2.0 and 2.0 );
\fill [line width=1.0pt, custcol_0_0_0]  (785.0,315.0) ellipse (2.0 and 2.0 );
\fill [line width=1.0pt, custcol_0_0_0]  (785.0,335.0) ellipse (2.0 and 2.0 );
\fill [line width=1.0pt, custcol_0_0_0]  (785.0,355.0) ellipse (2.0 and 2.0 );
\fill [line width=1.0pt, custcol_0_0_0]  (805.0,365.0) ellipse (2.0 and 2.0 );
\fill [line width=1.0pt, custcol_0_0_0]  (825.0,325.0) ellipse (2.0 and 2.0 );
\draw [line width=3.0pt, custcol_0_0_0 ]  (885.0,45.0) -- (965.0,45.0) -- (965.0,55.0) -- (985.0,55.0) -- (987.0,55.0) ;
\draw [line width=3.0pt, custcol_0_0_0 ]  (955.0,85.0) -- (965.0,85.0) -- (965.0,75.0) -- (985.0,75.0) -- (987.0,75.0) ;
\draw [line width=2.0pt, custcol_0_0_0 ]  (1015.0,65.0) .. controls  (1005.0,50.0)  ..  (985.0,50.0) .. controls  (993.0,65.0)  ..  (985.0,80.0) .. controls  (1005.0,80.0)  ..  (1015.0,65.0) -- cycle ;
\logisimfontC{\fontsize{10pt}{10pt}\fontseries{bx}\selectfont\node[inner sep=0, outer sep=0, custcol_0_0_0, anchor=base west] at  (29.0,37.0)  {PCctrl};}
\draw [line width=2.0pt, custcol_0_0_0 ]  (26.0,24.0) -- (65.0,24.0) -- (65.0,42.0) -- (55.0,42.0) -- (45.0,46.0) -- (35.0,42.0) -- (26.0,42.0) -- cycle;
\fill [line width=2.0pt, custcol_0_0_0]  (45.0,45.0) ellipse (2.0 and 2.0 );
\logisimfontC{\fontsize{10pt}{10pt}\fontseries{bx}\selectfont\node[inner sep=0, outer sep=0, custcol_0_0_0, anchor=base west] at  (783.0,400.0)  {ALUsrcB};}
\draw [line width=2.0pt, custcol_0_0_0 ]  (780.0,389.0) -- (795.0,389.0) -- (805.0,385.0) -- (815.0,389.0) -- (831.0,389.0) -- (831.0,406.0) -- (780.0,406.0) -- cycle;
\fill [line width=2.0pt, custcol_0_0_0]  (805.0,385.0) ellipse (2.0 and 2.0 );
\logisimfontC{\fontsize{8pt}{8pt}\fontseries{bx}\selectfont\node[inner sep=0, outer sep=0, custcol_0_0_0, anchor=base west] at  (214.0,88.0)  {memWrite};}
\draw [line width=2.0pt, custcol_0_0_0 ]  (211.0,77.0) -- (260.0,77.0) -- (260.0,92.0) -- (245.0,92.0) -- (235.0,96.0) -- (225.0,92.0) -- (211.0,92.0) -- cycle;
\fill [line width=2.0pt, custcol_0_0_0]  (235.0,95.0) ellipse (2.0 and 2.0 );
\draw [line width=4.0pt, custcol_0_0_0]  (625.0,405.0) ellipse (49.0 and 50.0 );
\logisimfontC{\fontsize{16pt}{16pt}\fontseries{bx}\selectfont\node[inner sep=0, outer sep=0, custcol_0_0_0, anchor=base west] at  (581.0,404.0)  {Immediate};}
\logisimfontC{\fontsize{16pt}{16pt}\fontseries{bx}\selectfont\node[inner sep=0, outer sep=0, custcol_0_0_0, anchor=base west] at  (582.0,417.0)  {Generator};}
\fill [line width=1.0pt, custcol_0_0_0]  (575.0,405.0) ellipse (2.0 and 2.0 );
\fill [line width=1.0pt, custcol_0_0_0]  (675.0,405.0) ellipse (2.0 and 2.0 );
\logisimfontC{\fontsize{10pt}{10pt}\fontseries{bx}\selectfont\node[inner sep=0, outer sep=0, custcol_0_0_0, anchor=base west] at  (839.0,49.0)  {PCwrite};}
\draw [line width=2.0pt, custcol_0_0_0 ]  (836.0,36.0) -- (882.0,36.0) -- (886.0,45.0) -- (882.0,55.0) -- (836.0,55.0) -- cycle;
\fill [line width=2.0pt, custcol_0_0_0]  (885.0,45.0) ellipse (2.0 and 2.0 );
\logisimfontB{\fontsize{14pt}{14pt}\selectfont\node[inner sep=0, outer sep=0, custcol_87_87_87, anchor=base west] at  (791.0,160.0)  {0};}
\draw [line width=4.0pt, custcol_0_0_0 ]  (785.0,135.0) -- (785.0,235.0) -- (825.0,215.0) -- (825.0,155.0) -- cycle;
\logisimfontC{\fontsize{14pt}{14pt}\fontseries{bx}\selectfont\node[inner sep=0, outer sep=0, custcol_0_0_0, anchor=base west] at  (791.0,185.0)  {Mux};}
\logisimfontC{\fontsize{14pt}{14pt}\fontseries{bx}\selectfont\node[inner sep=0, outer sep=0, custcol_0_0_0, anchor=base west] at  (794.0,196.0)  {4x1};}
\fill [line width=1.0pt, custcol_0_0_0]  (785.0,155.0) ellipse (2.0 and 2.0 );
\fill [line width=1.0pt, custcol_0_0_0]  (785.0,175.0) ellipse (2.0 and 2.0 );
\fill [line width=1.0pt, custcol_0_0_0]  (785.0,195.0) ellipse (2.0 and 2.0 );
\fill [line width=1.0pt, custcol_0_0_0]  (785.0,215.0) ellipse (2.0 and 2.0 );
\fill [line width=1.0pt, custcol_0_0_0]  (805.0,225.0) ellipse (2.0 and 2.0 );
\fill [line width=1.0pt, custcol_0_0_0]  (825.0,185.0) ellipse (2.0 and 2.0 );
\draw [line width=4.0pt, custcol_0_0_0 ]  (25.0,55.0) -- (94.0,55.0) ;
\draw [line width=4.0pt, custcol_0_0_0 ]  (95.0,55.0) -- (95.0,174.0) ;
\draw [line width=4.0pt, custcol_0_0_0 ]  (95.0,175.0) -- (26.0,175.0) ;
\draw [line width=4.0pt, custcol_0_0_0 ]  (25.0,175.0) -- (25.0,56.0) ;
\draw [line width=1.0pt, custcol_0_0_0 ]  (69.0,175.0) -- (75.0,165.0) -- (81.0,175.0) ;
\logisimfontC{\fontsize{20pt}{20pt}\fontseries{bx}\selectfont\node[inner sep=0, outer sep=0, custcol_0_0_0, anchor=base west] at  (45.0,123.0)  {PC};}
\draw [line width=2.0pt, custcol_0_0_0 ]  (95.0,60.0) -- (101.0,60.0) -- (101.0,181.0) -- (30.0,181.0) -- (30.0,175.0) ;
\draw [line width=2.0pt, custcol_0_0_0 ]  (101.0,65.0) -- (106.0,65.0) -- (106.0,186.0) -- (35.0,186.0) -- (35.0,182.0) ;
\fill [line width=1.0pt, custcol_0_0_0]  (25.0,115.0) ellipse (2.0 and 2.0 );
\fill [line width=1.0pt, custcol_0_0_0]  (45.0,55.0) ellipse (2.0 and 2.0 );
\fill [line width=1.0pt, custcol_0_0_0]  (75.0,175.0) ellipse (2.0 and 2.0 );
\fill [line width=1.0pt, custcol_0_0_0]  (95.0,115.0) ellipse (2.0 and 2.0 );
\logisimfontC{\fontsize{10pt}{10pt}\fontseries{bx}\selectfont\node[inner sep=0, outer sep=0, custcol_0_0_0, anchor=base west] at  (843.0,79.0)  {PCwriteCond};}
\draw [line width=2.0pt, custcol_0_0_0 ]  (840.0,66.0) -- (912.0,66.0) -- (916.0,75.0) -- (912.0,85.0) -- (840.0,85.0) -- cycle;
\fill [line width=2.0pt, custcol_0_0_0]  (915.0,75.0) ellipse (2.0 and 2.0 );
\draw [line width=1.0pt, custcol_0_0_0 ]  (666.0,49.0) -- (669.0,49.0) -- (669.0,89.0) -- (620.0,89.0) -- (620.0,86.0) ;
\draw [line width=1.0pt, custcol_0_0_0 ]  (653.0,85.0) -- (655.0,78.0) -- (657.0,85.0) ;
\draw [line width=2.0pt, custcol_0_0_0 ]  (615.0,45.0) -- (664.0,45.0) ;
\draw [line width=2.0pt, custcol_0_0_0 ]  (665.0,45.0) -- (665.0,84.0) ;
\draw [line width=2.0pt, custcol_0_0_0 ]  (665.0,85.0) -- (616.0,85.0) ;
\draw [line width=2.0pt, custcol_0_0_0 ]  (615.0,85.0) -- (615.0,46.0) ;
\draw [line width=1.0pt, custcol_0_0_0 ]  (670.0,52.0) -- (672.0,52.0) -- (672.0,92.0) -- (625.0,92.0) -- (625.0,90.0) ;
\logisimfontC{\fontsize{11pt}{11pt}\fontseries{bx}\selectfont\node[inner sep=0, outer sep=0, custcol_0_0_0, anchor=base west] at  (618.0,69.0)  {PCback};}
\fill [line width=1.0pt, custcol_0_0_0]  (615.0,65.0) ellipse (2.0 and 2.0 );
\fill [line width=1.0pt, custcol_0_0_0]  (625.0,45.0) ellipse (2.0 and 2.0 );
\fill [line width=1.0pt, custcol_0_0_0]  (655.0,85.0) ellipse (2.0 and 2.0 );
\fill [line width=1.0pt, custcol_0_0_0]  (665.0,65.0) ellipse (2.0 and 2.0 );
\logisimfontC{\fontsize{16pt}{16pt}\fontseries{bx}\selectfont\node[inner sep=0, outer sep=0, custcol_0_0_0, anchor=base west] at  (760.0,320.0)  {4};}
\logisimfontC{\fontsize{9pt}{9pt}\fontseries{bx}\selectfont\node[inner sep=0, outer sep=0, custcol_0_0_0, anchor=base west] at  (473.0,459.0)  {mem2Reg};}
\draw [line width=2.0pt, custcol_0_0_0 ]  (470.0,449.0) -- (485.0,449.0) -- (495.0,445.0) -- (505.0,449.0) -- (521.0,449.0) -- (521.0,464.0) -- (470.0,464.0) -- cycle;
\fill [line width=2.0pt, custcol_0_0_0]  (495.0,445.0) ellipse (2.0 and 2.0 );
\draw [line width=1.0pt, custcol_0_0_0 ]  (409.0,505.0) -- (415.0,495.0) -- (421.0,505.0) ;
\draw [line width=2.0pt, custcol_0_0_0 ]  (443.0,395.0) -- (447.0,395.0) -- (447.0,517.0) -- (375.0,517.0) -- (375.0,513.0) ;
\draw [line width=4.0pt, custcol_0_0_0 ]  (365.0,385.0) -- (434.0,385.0) ;
\draw [line width=4.0pt, custcol_0_0_0 ]  (435.0,385.0) -- (435.0,504.0) ;
\draw [line width=4.0pt, custcol_0_0_0 ]  (435.0,505.0) -- (366.0,505.0) ;
\draw [line width=4.0pt, custcol_0_0_0 ]  (365.0,505.0) -- (365.0,386.0) ;
\draw [line width=2.0pt, custcol_0_0_0 ]  (435.0,390.0) -- (442.0,390.0) -- (442.0,512.0) -- (370.0,512.0) -- (370.0,505.0) ;
\logisimfontC{\fontsize{11pt}{11pt}\fontseries{bx}\selectfont\node[inner sep=0, outer sep=0, custcol_0_0_0, anchor=base west] at  (386.0,444.0)  {Data};}
\logisimfontC{\fontsize{11pt}{11pt}\fontseries{bx}\selectfont\node[inner sep=0, outer sep=0, custcol_0_0_0, anchor=base west] at  (375.0,454.0)  {Register};}
\fill [line width=1.0pt, custcol_0_0_0]  (365.0,445.0) ellipse (2.0 and 2.0 );
\fill [line width=1.0pt, custcol_0_0_0]  (415.0,505.0) ellipse (2.0 and 2.0 );
\fill [line width=1.0pt, custcol_0_0_0]  (435.0,445.0) ellipse (2.0 and 2.0 );
\logisimfontC{\fontsize{10pt}{10pt}\fontseries{bx}\selectfont\node[inner sep=0, outer sep=0, custcol_0_0_0, anchor=base west] at  (592.0,27.0)  {PCbackWrite};}
\draw [line width=2.0pt, custcol_0_0_0 ]  (589.0,14.0) -- (662.0,14.0) -- (662.0,32.0) -- (635.0,32.0) -- (625.0,36.0) -- (615.0,32.0) -- (589.0,32.0) -- cycle;
\fill [line width=2.0pt, custcol_0_0_0]  (625.0,35.0) ellipse (2.0 and 2.0 );
\logisimfontC{\fontsize{10pt}{10pt}\fontseries{bx}\selectfont\node[inner sep=0, outer sep=0, custcol_0_0_0, anchor=base west] at  (366.0,137.0)  {IRWrite};}
\draw [line width=2.0pt, custcol_0_0_0 ]  (363.0,124.0) -- (407.0,124.0) -- (407.0,142.0) -- (395.0,142.0) -- (385.0,146.0) -- (375.0,142.0) -- (363.0,142.0) -- cycle;
\fill [line width=2.0pt, custcol_0_0_0]  (385.0,145.0) ellipse (2.0 and 2.0 );
\draw [line width=1.0pt, custcol_0_0_0 ]  (736.0,279.0) -- (739.0,279.0) -- (739.0,319.0) -- (699.0,319.0) -- (699.0,316.0) ;
\logisimfontC{\fontsize{22pt}{22pt}\fontseries{bx}\selectfont\node[inner sep=0, outer sep=0, custcol_0_0_0, anchor=base west] at  (707.0,304.0)  {B};}
\draw [line width=1.0pt, custcol_0_0_0 ]  (739.0,282.0) -- (742.0,282.0) -- (742.0,322.0) -- (702.0,322.0) -- (702.0,320.0) ;
\draw [line width=2.0pt, custcol_0_0_0 ]  (695.0,275.0) -- (734.0,275.0) ;
\draw [line width=2.0pt, custcol_0_0_0 ]  (735.0,275.0) -- (735.0,314.0) ;
\draw [line width=2.0pt, custcol_0_0_0 ]  (735.0,315.0) -- (696.0,315.0) ;
\draw [line width=2.0pt, custcol_0_0_0 ]  (695.0,315.0) -- (695.0,276.0) ;
\draw [line width=1.0pt, custcol_0_0_0 ]  (723.0,315.0) -- (725.0,308.0) -- (727.0,315.0) ;
\fill [line width=1.0pt, custcol_0_0_0]  (695.0,295.0) ellipse (2.0 and 2.0 );
\fill [line width=1.0pt, custcol_0_0_0]  (725.0,315.0) ellipse (2.0 and 2.0 );
\fill [line width=1.0pt, custcol_0_0_0]  (735.0,295.0) ellipse (2.0 and 2.0 );
\draw [line width=3.0pt, custcol_0_0_0 ]  (765.0,375.0) -- (765.0,380.0) ;
\draw [line width=1.0pt, custcol_0_0_0 ]  (773.0,381.0) -- (757.0,381.0) ;
\draw [line width=1.0pt, custcol_0_0_0 ]  (770.0,384.0) -- (760.0,384.0) ;
\draw [line width=1.0pt, custcol_0_0_0 ]  (767.0,387.0) -- (763.0,387.0) ;
\fill [line width=1.0pt, custcol_0_0_0]  (765.0,375.0) ellipse (2.0 and 2.0 );
\end{tikzpicture}
}
    \vspace{-15pt}
    \caption{Esquemático do Multiciclo com ISA reduzida}
    \vspace{-5pt}
\end{figure}

\noindent Em que os túneis representam entradas e saídas do módulo de controle do processador, que foi abstraído na figura acima e será devidamente explicado na seção adequada.

%\begin{figure}[H]
%    \centering
%    %\includegraphics{Recursos/Imagens/CONTROL.png}
%    \caption{Módulo de controle do Multiciclo %com ISA reduzida}
%\end{figure}

\subsubsection{Tarefa I}
\begin{tcolorbox}[title=Enunciado, colback=blue!5!white, colframe=blue!75!black]

\end{tcolorbox}


\subsubsection{Tarefa II}
\begin{tcolorbox}[title=Enunciado, colback=blue!5!white, colframe=blue!75!black]


\end{tcolorbox}

\subsubsection{Tarefa III}
\begin{tcolorbox}[title=Enunciado, colback=blue!5!white, colframe=blue!75!black]

Implemente o Bloco Controlador e desenhe a máquina de estados do controle. 
\end{tcolorbox}
Como o enunciado adianta, o bloco de controle do multiciclo é implementado via máquina de estado. Particularmente, utilizamos uma máquina de estados finita do tipo Moore (isto é, em que a saída depende exclusivamente do estado atual, não das entradas). Para projetar uma máquina de estados, primeiro precisamos do seu diagrama de estados; o diagrama de estados completo do bloco de controle do multiciclo com ISA reduzida, considerando as doze instruções solicitadas, é mostrado na \autoref{fig:diag_estados}.

\begin{figure}[H]
    \centering
    \resizebox{\textwidth}{!}{
\definecolor{amarelo_queimado}{HTML}{FFC222}
\begin{tikzpicture}[
    % ========== estilos ==========
    state/.style = {
        rectangle split,
        rectangle split parts=2,
        rectangle split horizontal=false,
        draw=black,
        very thick,
        rounded corners,
        rectangle split part fill={amarelo_queimado,white},
        text width=50mm,
        align=left,
        inner sep=6pt,
        font=\footnotesize
    },
    bitname/.style = {anchor=west, font=\footnotesize},
    bitbar/.style = {anchor=west, font=\footnotesize, text=black},
    arrow/.style = {-{Latex[length=3mm,width=2.5mm]}, semithick}
]

% ======= Nodes (layout em camadas para evitar cruzamentos) =======
% Top layer

\node[state] (Fetch) at (0,10) {
    \centering \normalsize \textbf{Fetch}
\nodepart{two}
\parbox{50mm}{\centering
ALUsrcA = 10, ALUsrcB = 01, ALUOp = 00,
$\overline{\mathrm{PCsrc}}$, PCbackWrite, PCWrite,
$\overline{\mathrm{IorD}}$, MemRead, IRWrite}
};

\node[state] (Decode) at (-3.5,6) {
    \centering \normalsize \textbf{Decode}
\nodepart{two}
\parbox{50mm}{\centering
ALUsrcA = 00,
ALUsrcB = 10,
ALUOp = 00}
};

\node[state] (WBLui) at (3.5,6) {
    \centering \normalsize \textbf{WBLui}
\nodepart{two}
\parbox{50mm}{\centering
Mem2Reg = 11, RegWrite}
};

% Middle layer: left cluster (Imm path), center (WBLui), right cluster (RegOp & BrCheck & Jal)
\node[state] (ImmSum) at (-10,3) {
    \centering \normalsize \textbf{ImmSum}
\nodepart{two}
\parbox{50mm}{\centering
ALUsrcA = 01,
ALUsrcB = 10,
ALUOp = 00}
};

\node[state] (RegOp) at (-3.25,3) {
    \centering \normalsize \textbf{RegOp}
\nodepart{two}
\parbox{50mm}{\centering
ALUsrcA = 01,
ALUsrcB = 00,
ALUOp = 10}
};

\node[state] (BrCheck) at (3.25,3) {
    \centering \normalsize \textbf{BrCheck}
\nodepart{two}
\parbox{50mm}{\centering
ALUOp = 01, ALUsrcA = 01, ALUsrcB = 00, PCsrc, PCwriteCond}
};

\node[state] (Jal) at (10,3) {
    \centering \normalsize \textbf{Jal}
\nodepart{two}
\parbox{50mm}{\centering
Mem2Reg = 01, RegWrite, PCsrc, PCWrite}
};

% Bottom layer: memory and writeback nodes
\node[state] (MemRead) at (-10,0) {
    \centering \normalsize \textbf{MemRead}
\nodepart{two}
\parbox{50mm}{\centering
IorD, MemRead}
};

\node[state] (MemWrite) at (-3,0) {
    \centering \normalsize \textbf{MemWrite}
\nodepart{two}
\parbox{50mm}{\centering
IorD, MemWrite}
};

\node[state] (WB_ALU) at (4.5,0) {
    \centering \normalsize \textbf{WB\_ALU}
\nodepart{two}
\parbox{50mm}{\centering
Mem2Reg = 00, RegWrite}
};

\node[state] (WBMem) at (-6.5,-3) {
    \centering \normalsize \textbf{WBMem}
\nodepart{two}
\parbox{50mm}{\centering
RegWrite, Mem2Reg = 10
}
};

% ======= Edges (curvas com out/in para evitar nós) =======
% Fetch -> Decode
\draw[arrow] ($(Fetch.south west) + (12mm, 0mm)$) to[out=-120, in=80] node[midway,above,sloped] {\scriptsize 1} (Decode.north);

% Decode -> ImmSum (LW/SW/BRANCH)
\draw[arrow] ($(Decode.west) + (0mm, -3mm)$) to[out=180,in=90] node[pos=0.4, above, yshift=10pt, sloped] {\scriptsize 0000011, 0100011,} node[pos=0.4, above, sloped]{\scriptsize 0010011, 1100111} (ImmSum.north);

% Decode -> RegOp (R-type)
\draw[arrow] (Decode.south) -- node[midway, left] {\scriptsize 0110011} (RegOp.north);

% Decode -> BrCheck (BEQ) (route passing to right)
\draw[arrow] ($(Decode.south east) + (-0.7mm, 0.7mm)$) -- node[pos=0.35, below, sloped] {\scriptsize 1100011} (BrCheck.north west);

% Decode -> Jal
\draw[arrow] ($(Decode.south east) + (0,5mm)$) to[out=-30,in=145] node[midway, below, sloped] {\scriptsize 1101111} (Jal.north west);

% ImmSum -> MemRead (LW)
\draw[arrow] (ImmSum.south) to[out=-90,in=90] node[pos=0.3, left] {\scriptsize 0000011} (MemRead.north);

% ImmSum -> MemWrite (SW) (route under center)
\draw[arrow,rounded corners=6pt,shorten >=1pt,shorten <=1pt] 
    ($(ImmSum.south east) + (-20mm,0)$)
    to[out=-90,in=150,looseness=1.2] node[pos=0.6, above, sloped] {\scriptsize 0100011} ($(MemRead.north east)+(2mm, 1mm)$)
    to[out=-60, in=170, looseness=1.0] (MemWrite.west);

% ImmSum -> WB_ALU (addi)
\draw[arrow,rounded corners=6pt,shorten >=1pt,shorten <=1pt] 
    ($(ImmSum.south east) + (-5mm,0)$)
    to[out=-80,in=170,looseness=1.0] ($(MemRead.north east)+(6mm, 5mm)$)
    to[out=-5,in=160,looseness=1.0] node[pos=0.3, above, sloped] {\scriptsize 0010011} ($(MemWrite.north east)+(2mm, 2mm)$)
    -- ($(WB_ALU.west)+(0,3mm)$);

% MemRead -> WBMem -> Fetch
\draw[arrow] (MemRead.south) to[out=-90,in=90] node[midway, above, sloped] {\scriptsize 1} (WBMem.north);
\draw[arrow,rounded corners=6pt,shorten >=1pt,shorten <=1pt]
    (WBMem.west) -- ($(WBMem.west)!0.4!(WBMem.west -| ImmSum.west)$)
    to[out=170,in=-90,looseness=1.1] ($(ImmSum.west)+(-10mm,0mm)$)
    to[out=90,in=180,looseness=1.5] node[pos=0, left] {\scriptsize 1} ($(Fetch.west)+(-40mm,0mm)$)
    -- (Fetch.west);

% MemWrite -> Fetch
\draw[-,rounded corners=6pt,shorten >=1pt,shorten <=1pt]
    ($(MemWrite.south east) + (-15mm,0)$)
    to[out=-30,in=180,looseness=1] ($(WBMem.center -| Jal.south) + (-20mm, 0)$)
    to[out=0, in=-100, looseness=1.2] ($(Jal.south east) + (9mm,-2mm)$);

% RegOp -> WB_ALU -> Fetch
\draw[arrow,rounded corners=6pt,shorten >=1pt,shorten <=1pt]
    ($(RegOp.south east) + (-1mm,1mm)$) -- node[pos=0.3, above, sloped] {\scriptsize 1} (WB_ALU.north west);
\draw[arrow,rounded corners=6pt,shorten >=1pt,shorten <=1pt]
    (WB_ALU.east) -- ($(WB_ALU.east)!0.4!(WB_ALU -| Jal)$)
    to[out=0,in=-90, looseness=1.3] ($(Jal.east)+(10mm,0)$)
    to[out=90,in=0,looseness=1.5] node[pos=0.2, right] {\scriptsize 1} ($(Fetch.east)+(40mm,0)$)
    -- (Fetch.east);

% BrCheck -> Fetch
\draw[-] ($(BrCheck.south east) + (0,5mm)$) to[out=-45,in=175] ($(WB_ALU.east)+(25mm,0)$);

% Jal -> Fetch
\draw[arrow] (Jal.north) to[out=90,in=0] node[midway, above, sloped] {\scriptsize 1} ($(Fetch.south east) + (0,5mm)$);

% Decode -> WBLui
\draw[arrow] (Decode.east) -- node[pos=0.4, above, sloped] {\scriptsize 0110111} (WBLui.west);
\draw[arrow] (WBLui.north) to[out=120,in=-60] node[midway, above, sloped] {\scriptsize 1} ($(Fetch.south east) + (-15mm,0)$);

\end{tikzpicture}
}
    \vspace{-15pt}
    \caption{Diagrama de estados do controle}
    \label{fig:diag_estados}
\end{figure}

A partir desse diagrama, conseguimos montar as tabelas \hyperref[tab:saidas]{(1)} de saída e \hyperref[tab:tran_estados]{(2)} de transição de estados, a partir das quais é possível inferir as lógicas de saída e de próximo estado, respectivamente.


\begin{table}[H]
    \centering
    \resizebox{\textwidth}{!}{
    \begin{tabular}{|c|c|c|c|c|c|c|c|c|c|c|c|c|c|}
        \hline
        \rowcolor{black}
        \textcolor{white}{\textbf{Estado}} & \multicolumn{13}{|c|}{\textcolor{white}{\textbf{Saídas}}} \\ \cline{2-14}
        
        \rowcolor{black}
        \textcolor{white}{\textbf{Atual}} &
        \textcolor{white}{\textbf{\scriptsize PCsrc}} &
        \textcolor{white}{\textbf{\scriptsize ALUOp}} &
        \textcolor{white}{\textbf{\scriptsize ALUsrcA}} &
        \textcolor{white}{\textbf{\scriptsize ALUsrcB}} &
        \textcolor{white}{\textbf{\scriptsize PCbackWrite}} &
        \textcolor{white}{\textbf{\scriptsize RegWrite}} &
        \textcolor{white}{\textbf{\scriptsize Mem2Reg}} &
        \textcolor{white}{\textbf{\scriptsize PCwriteCond}} &
        \textcolor{white}{\textbf{\scriptsize PCWrite}} &
        \textcolor{white}{\textbf{\scriptsize IorD}} &
        \textcolor{white}{\textbf{\scriptsize MemWrite}} &
        \textcolor{white}{\textbf{\scriptsize MemRead}} &
        \textcolor{white}{\textbf{\scriptsize IRWrite}}
        \\ \hline
         
        0000 & 0 & 00 & 10 & 01 & 1 & 0 & XX & 0 & 1 & 0 & 0 & 1 & 1 \\ \hline

        \rowcolor{gray!20}
        0001 & X & 00 & 00 & 10 & 0 & 0 & XX & 0 & 0 & X & 0 & 0 & 0 \\ \hline
        
        0011 & X & 00 & 01 & 10 & 0 & 0 & XX & 0 & 0 & X & 0 & 0 & 0 \\ \hline

        \rowcolor{gray!20}
        0010 & X & 10 & 01 & 00 & 0 & 0 & XX & 0 & 0 & X & 0 & 0 & 0 \\ \hline

        1000 & 1 & 01 & 01 & 00 & 0 & 0 & XX & 1 & 0 & X & 0 & 0 & 0 \\ \hline

        \rowcolor{gray!20}
        0101 & 1 & XX & XX & XX & 0 & 1 & 01 & 0 & 1 & X & 0 & 0 & 0 \\ \hline

        0110 & X & XX & XX & XX & 0 & 0 & XX & 0 & 0 & 1 & 0 & 1 & 0 \\ \hline

        \rowcolor{gray!20}
        0111 & X & XX & XX & XX & 0 & 0 & XX & 0 & 0 & 1 & 1 & 0 & 0 \\ \hline

        1001 & X & XX & XX & XX & 0 & 1 & 00 & 0 & 0 & X & 0 & 0 & 0 \\ \hline

        \rowcolor{gray!20}
        1011 & X & XX & XX & XX & 0 & 1 & 10 & 0 & 0 & X & 0 & 0 & 0 \\ \hline

        1101 & 0 & 00 & 10 & 01 & 0 & 1 & 11 & 0 & 0 & X & 0 & 0 & 0 \\ \hline
    \end{tabular}
    }
    \caption{Tabela verdade das saídas do controle}
    \label{tab:saidas}
\end{table}

Com base nesta tabela, utilizamos Mapas de Karnaugh para minimizar, na forma de soma de produtos, as expressões para as saídas. Os resultados obtidos estão exibidos a seguir, em que $Q_i$ representa o $i$-ésimo bit do estado atual.
{\footnotesize
\begin{align*}
    &PCsrc =
        \overline{Q_3}\,Q_2
        +
        Q_3\,\overline{Q_2}\,\overline{Q_1}\\
    &ALUOp_1 =
        \overline{Q_3}\,Q_1\,\overline{Q_0}, \
    ALUOp_0 =
        Q_3\,\overline{Q_2}\,\overline{Q_1}\\
    &ALUsrcA_1 =
        Q_2\,\overline{Q_1}\,Q_0
        +
        \overline{Q_3}\,\overline{Q_2}\,\overline{Q_1}\,\overline{Q_0}, \
    ALUsrcA_0 =
        \overline{Q_3}\,Q_1
        +
        Q_3\,\overline{Q_2}\,\overline{Q_1}\\
    &ALUsrcB_1 =
        \overline{Q_3}\,Q_0, \
    ALUsrcB_0 =
        Q_2\,\overline{Q_1}\,Q_0
        +
        \overline{Q_3}\,\overline{Q_2}\,\overline{Q_1}\,\overline{Q_0}\\
    &PCbackWrite =
        \overline{Q_3}\,\overline{Q_2}\,\overline{Q_1}\,\overline{Q_0}\\
    &RegWrite =
        Q_3\,\overline{Q_2}\,Q_0
        +
        Q_2\,\overline{Q_1}\,Q_0\\
    &Mem2Reg_1 =
        \overline{Q_2}\,Q_1\,Q_0
        +
        Q_3\,Q_2\,\overline{Q_1}\,Q_0, \
    Mem2Reg_0 =
        Q_2\,\overline{Q_1}\,Q_0\\
    &PCWriteCond = 
    Q_3\,\overline{Q_2}\,\overline{Q_1}\,\overline{Q_0}\\
    &PCWrite =
        \overline{Q_3}\,\overline{Q_2}\,\overline{Q_1}\,\overline{Q_0}
        +
        \overline{Q_3}\,Q_2\,\overline{Q_1}\,Q_0 \\
    &IorD =
        \overline{Q_3}\,Q_1\\
    &MemWrite =
        \overline{Q_3}\,Q_2\,Q_1\,Q_0\\
    &MemRead =
        \overline{Q_3}\,\overline{Q_2}\,\overline{Q_1}\,\overline{Q_0}
        +
        \overline{Q_3}\,Q_2\,Q_1\,\overline{Q_0}\\
    &IRWrite =
        \overline{Q_3}\,\overline{Q_2}\,\overline{Q_1}\,\overline{Q_0}
\end{align*}
}

\noindent Tendo essas equações, é fácil implementar a lógica de saída. O esquemático a seguir mostra essa implemetação.

\begin{figure}[H]
    \centering
    % Important: If latex complains about unicode characters,
% please use "\usepackage[utf8x]{inputenc}" in your preamble
% You can change the size of the picture by putting it into the construct:
% 1) \resizebox{10cm}{!}{"below picture"} to scale horizontally to 10 cm
% 2) \resizebox{!}{15cm}{"below picture"} to scale vertically to 15 cm
% 3) \resizebox{10cm}{15cm}{"below picture"} a combination of above two
% It is not recomended to use the scale option of the tikzpicture environment.
\resizebox{\textwidth}{!}{
\begin{tikzpicture}[x=1pt,y=-1pt,line cap=rect]
\def\logisimfontA#1{\fontfamily{cmr}{#1}} % Replaced by logisim, original font was "SansSerif"
\def\logisimfontB#1{\fontfamily{cmtt}{#1}} % Replaced by logisim, original font was "Monospaced"
\def\logisimfontC#1{\fontfamily{Arial Rounded MT Bold}{#1}}
\definecolor{custcol_0_0_0}{RGB}{0, 0, 0}
\definecolor{custcol_ff_ff_ff}{RGB}{255, 255, 255}
\draw [line width=3.0pt, custcol_0_0_0 ]  (765.0,185.0) -- (765.0,235.0) ;
\draw [line width=3.0pt, custcol_0_0_0 ]  (725.0,95.0) -- (725.0,65.0) -- (675.0,65.0) -- (675.0,95.0) ;
\draw [line width=3.0pt, custcol_0_0_0 ]  (1215.0,105.0) -- (1215.0,25.0) -- (1165.0,25.0) -- (1165.0,95.0) ;
\draw [line width=3.0pt, custcol_0_0_0 ]  (835.0,65.0) -- (835.0,105.0) ;
\draw [line width=3.0pt, custcol_0_0_0 ]  (965.0,65.0) -- (915.0,65.0) -- (915.0,105.0) ;
\draw [line width=4.0pt, custcol_0_0_0 ]  (475.0,225.0) -- (475.0,215.0) -- (495.0,215.0) ;
\draw [line width=3.0pt, custcol_0_0_0 ]  (1185.0,45.0) -- (1135.0,45.0) -- (1135.0,105.0) ;
\draw [line width=3.0pt, custcol_0_0_0 ]  (465.0,95.0) -- (465.0,45.0) -- (415.0,45.0) -- (415.0,95.0) ;
\draw [line width=3.0pt, custcol_0_0_0 ]  (1275.0,135.0) -- (1275.0,235.0) ;
\draw [line width=3.0pt, custcol_0_0_0 ]  (385.0,95.0) -- (385.0,45.0) -- (335.0,45.0) -- (335.0,105.0) ;
\draw [line width=3.0pt, custcol_0_0_0 ]  (1165.0,25.0) -- (1115.0,25.0) -- (1115.0,105.0) ;
\draw [line width=3.0pt, custcol_0_0_0 ]  (625.0,65.0) -- (625.0,105.0) ;
\draw [line width=3.0pt, custcol_0_0_0 ]  (895.0,105.0) -- (895.0,25.0) -- (935.0,25.0) -- (985.0,25.0) -- (1035.0,25.0) -- (1035.0,105.0) ;
\draw [line width=3.0pt, custcol_0_0_0 ]  (985.0,25.0) -- (985.0,95.0) ;
\draw [line width=3.0pt, custcol_0_0_0 ]  (1255.0,95.0) -- (1255.0,5.0) -- (1205.0,5.0) -- (1205.0,95.0) ;
\draw [line width=3.0pt, custcol_0_0_0 ]  (745.0,25.0) -- (745.0,95.0) ;
\draw [line width=3.0pt, custcol_0_0_0 ]  (565.0,5.0) -- (635.0,5.0) -- (635.0,95.0) ;
\draw [line width=3.0pt, custcol_0_0_0 ]  (1005.0,95.0) -- (1005.0,45.0) -- (955.0,45.0) -- (955.0,95.0) ;
\draw [line width=3.0pt, custcol_0_0_0 ]  (1015.0,185.0) -- (1015.0,235.0) ;
\draw [line width=3.0pt, custcol_0_0_0 ]  (295.0,25.0) -- (295.0,95.0) ;
\draw [line width=3.0pt, custcol_0_0_0 ]  (375.0,25.0) -- (375.0,95.0) ;
\draw [line width=3.0pt, custcol_0_0_0 ]  (875.0,95.0) -- (875.0,45.0) -- (905.0,45.0) -- (955.0,45.0) ;
\draw [line width=3.0pt, custcol_0_0_0 ]  (935.0,25.0) -- (935.0,95.0) ;
\draw [line width=3.0pt, custcol_0_0_0 ]  (1075.0,5.0) -- (1075.0,95.0) ;
\draw [line width=3.0pt, custcol_0_0_0 ]  (1205.0,5.0) -- (1155.0,5.0) -- (1155.0,95.0) ;
\draw [line width=3.0pt, custcol_0_0_0 ]  (695.0,25.0) -- (695.0,95.0) ;
\draw [line width=3.0pt, custcol_0_0_0 ]  (445.0,95.0) -- (445.0,25.0) -- (535.0,25.0) -- (535.0,95.0) ;
\draw [line width=3.0pt, custcol_0_0_0 ]  (755.0,105.0) -- (755.0,65.0) -- (795.0,65.0) -- (835.0,65.0) -- (885.0,65.0) -- (885.0,105.0) ;
\draw [line width=3.0pt, custcol_0_0_0 ]  (605.0,105.0) -- (605.0,25.0) -- (645.0,25.0) -- (695.0,25.0) -- (745.0,25.0) -- (775.0,25.0) -- (775.0,105.0) ;
\draw [line width=3.0pt, custcol_0_0_0 ]  (485.0,95.0) -- (485.0,5.0) -- (435.0,5.0) -- (435.0,95.0) ;
\draw [line width=3.0pt, custcol_0_0_0 ]  (1245.0,95.0) -- (1245.0,65.0) -- (1295.0,65.0) -- (1295.0,95.0) ;
\draw [line width=3.0pt, custcol_0_0_0 ]  (1205.0,185.0) -- (1205.0,235.0) ;
\draw [line width=3.0pt, custcol_0_0_0 ]  (715.0,45.0) -- (785.0,45.0) ;
\draw [line width=3.0pt, custcol_0_0_0 ]  (825.0,45.0) -- (825.0,105.0) ;
\draw [line width=3.0pt, custcol_0_0_0 ]  (535.0,25.0) -- (605.0,25.0) ;
\draw [line width=3.0pt, custcol_0_0_0 ]  (1125.0,135.0) -- (1125.0,235.0) ;
\draw [line width=3.0pt, custcol_0_0_0 ]  (715.0,95.0) -- (715.0,45.0) -- (665.0,45.0) -- (665.0,95.0) ;
\draw [line width=3.0pt, custcol_0_0_0 ]  (465.0,45.0) -- (505.0,45.0) -- (505.0,105.0) ;
\draw [line width=3.0pt, custcol_0_0_0 ]  (905.0,45.0) -- (905.0,95.0) ;
\draw [line width=3.0pt, custcol_0_0_0 ]  (545.0,95.0) -- (545.0,45.0) -- (615.0,45.0) -- (665.0,45.0) ;
\draw [line width=3.0pt, custcol_0_0_0 ]  (1105.0,5.0) -- (1105.0,95.0) ;
\draw [line width=3.0pt, custcol_0_0_0 ]  (645.0,25.0) -- (645.0,95.0) ;
\draw [line width=3.0pt, custcol_0_0_0 ]  (1025.0,5.0) -- (1025.0,95.0) ;
\draw [line width=3.0pt, custcol_0_0_0 ]  (265.0,105.0) -- (265.0,25.0) -- (295.0,25.0) -- (375.0,25.0) -- (405.0,25.0) -- (405.0,105.0) ;
\draw [line width=3.0pt, custcol_0_0_0 ]  (675.0,65.0) -- (625.0,65.0) -- (585.0,65.0) -- (475.0,65.0) -- (475.0,95.0) ;
\draw [line width=3.0pt, custcol_0_0_0 ]  (245.0,5.0) -- (285.0,5.0) -- (325.0,5.0) -- (325.0,95.0) ;
\draw [line width=3.0pt, custcol_0_0_0 ]  (1195.0,95.0) -- (1195.0,65.0) -- (1245.0,65.0) ;
\draw [line width=3.0pt, custcol_0_0_0 ]  (485.0,5.0) -- (525.0,5.0) -- (565.0,5.0) -- (565.0,95.0) ;
\draw [line width=3.0pt, custcol_0_0_0 ]  (795.0,65.0) -- (795.0,105.0) ;
\draw [line width=4.0pt, custcol_0_0_0 ]  (595.0,225.0) -- (595.0,215.0) -- (615.0,215.0) ;
\draw [line width=4.0pt, custcol_0_0_0 ]  (355.0,225.0) -- (355.0,215.0) -- (375.0,215.0) ;
\draw [line width=3.0pt, custcol_0_0_0 ]  (785.0,95.0) -- (785.0,45.0) -- (825.0,45.0) -- (855.0,45.0) -- (855.0,105.0) ;
\draw [line width=3.0pt, custcol_0_0_0 ]  (275.0,185.0) -- (275.0,235.0) ;
\draw [line width=3.0pt, custcol_0_0_0 ]  (615.0,45.0) -- (615.0,95.0) ;
\draw [line width=3.0pt, custcol_0_0_0 ]  (815.0,95.0) -- (815.0,25.0) -- (895.0,25.0) ;
\draw [line width=3.0pt, custcol_0_0_0 ]  (1055.0,95.0) -- (1055.0,45.0) -- (1015.0,45.0) -- (1015.0,95.0) ;
\draw [line width=3.0pt, custcol_0_0_0 ]  (775.0,25.0) -- (815.0,25.0) ;
\draw [line width=3.0pt, custcol_0_0_0 ]  (975.0,5.0) -- (975.0,95.0) ;
\draw [line width=3.0pt, custcol_0_0_0 ]  (345.0,65.0) -- (345.0,95.0) ;
\draw [line width=3.0pt, custcol_0_0_0 ]  (585.0,65.0) -- (585.0,105.0) ;
\draw [line width=3.0pt, custcol_0_0_0 ]  (855.0,45.0) -- (875.0,45.0) ;
\draw [line width=4.0pt, custcol_0_0_0 ]  (865.0,225.0) -- (865.0,215.0) -- (885.0,215.0) ;
\draw [line width=3.0pt, custcol_0_0_0 ]  (305.0,45.0) -- (335.0,45.0) ;
\draw [line width=3.0pt, custcol_0_0_0 ]  (685.0,235.0) -- (685.0,195.0) -- (705.0,195.0) -- (705.0,135.0) ;
\draw [line width=3.0pt, custcol_0_0_0 ]  (385.0,45.0) -- (415.0,45.0) ;
\draw [line width=3.0pt, custcol_0_0_0 ]  (725.0,65.0) -- (755.0,65.0) ;
\draw [line width=3.0pt, custcol_0_0_0 ]  (945.0,135.0) -- (945.0,235.0) ;
\draw [line width=3.0pt, custcol_0_0_0 ]  (425.0,65.0) -- (425.0,105.0) ;
\draw [line width=3.0pt, custcol_0_0_0 ]  (735.0,5.0) -- (735.0,105.0) ;
\draw [line width=3.0pt, custcol_0_0_0 ]  (885.0,65.0) -- (915.0,65.0) ;
\draw [line width=3.0pt, custcol_0_0_0 ]  (1065.0,65.0) -- (1065.0,105.0) ;
\draw [line width=3.0pt, custcol_0_0_0 ]  (1005.0,45.0) -- (1015.0,45.0) ;
\draw [line width=3.0pt, custcol_0_0_0 ]  (1195.0,65.0) -- (1145.0,65.0) -- (1145.0,105.0) ;
\draw [line width=3.0pt, custcol_0_0_0 ]  (635.0,5.0) -- (685.0,5.0) -- (735.0,5.0) -- (845.0,5.0) -- (925.0,5.0) -- (975.0,5.0) -- (1025.0,5.0) -- (1075.0,5.0) -- (1105.0,5.0) -- (1155.0,5.0) ;
\draw [line width=3.0pt, custcol_0_0_0 ]  (685.0,5.0) -- (685.0,95.0) ;
\draw [line width=3.0pt, custcol_0_0_0 ]  (1215.0,25.0) -- (1265.0,25.0) -- (1265.0,95.0) ;
\draw [line width=3.0pt, custcol_0_0_0 ]  (1185.0,95.0) -- (1185.0,45.0) -- (1235.0,45.0) -- (1285.0,45.0) -- (1285.0,95.0) ;
\draw [line width=3.0pt, custcol_0_0_0 ]  (405.0,25.0) -- (445.0,25.0) ;
\draw [line width=3.0pt, custcol_0_0_0 ]  (965.0,95.0) -- (965.0,65.0) -- (1065.0,65.0) -- (1095.0,65.0) -- (1145.0,65.0) ;
\draw [line width=3.0pt, custcol_0_0_0 ]  (505.0,45.0) -- (545.0,45.0) ;
\draw [line width=3.0pt, custcol_0_0_0 ]  (285.0,5.0) -- (285.0,105.0) ;
\draw [line width=3.0pt, custcol_0_0_0 ]  (1055.0,45.0) -- (1135.0,45.0) ;
\draw [line width=3.0pt, custcol_0_0_0 ]  (325.0,5.0) -- (365.0,5.0) -- (365.0,105.0) ;
\draw [line width=3.0pt, custcol_0_0_0 ]  (525.0,5.0) -- (525.0,105.0) ;
\draw [line width=3.0pt, custcol_0_0_0 ]  (845.0,5.0) -- (845.0,105.0) ;
\draw [line width=3.0pt, custcol_0_0_0 ]  (925.0,5.0) -- (925.0,105.0) ;
\draw [line width=3.0pt, custcol_0_0_0 ]  (1235.0,45.0) -- (1235.0,105.0) ;
\draw [line width=3.0pt, custcol_0_0_0 ]  (1095.0,65.0) -- (1095.0,105.0) ;
\draw [line width=3.0pt, custcol_0_0_0 ]  (1085.0,135.0) -- (1085.0,195.0) -- (1065.0,195.0) -- (1065.0,235.0) ;
\draw [line width=3.0pt, custcol_0_0_0 ]  (365.0,5.0) -- (435.0,5.0) ;
\draw [line width=3.0pt, custcol_0_0_0 ]  (1035.0,25.0) -- (1115.0,25.0) ;
\fill [line width=3.0pt, custcol_0_0_0]  (1055.0,45.0) ellipse (5.0 and 5.0 );
\fill [line width=3.0pt, custcol_0_0_0]  (625.0,65.0) ellipse (5.0 and 5.0 );
\fill [line width=3.0pt, custcol_0_0_0]  (1205.0,5.0) ellipse (5.0 and 5.0 );
\fill [line width=3.0pt, custcol_0_0_0]  (775.0,25.0) ellipse (5.0 and 5.0 );
\fill [line width=3.0pt, custcol_0_0_0]  (445.0,25.0) ellipse (5.0 and 5.0 );
\fill [line width=3.0pt, custcol_0_0_0]  (675.0,65.0) ellipse (5.0 and 5.0 );
\fill [line width=3.0pt, custcol_0_0_0]  (345.0,65.0) ellipse (5.0 and 5.0 );
\fill [line width=3.0pt, custcol_0_0_0]  (925.0,5.0) ellipse (5.0 and 5.0 );
\fill [line width=3.0pt, custcol_0_0_0]  (725.0,65.0) ellipse (5.0 and 5.0 );
\fill [line width=3.0pt, custcol_0_0_0]  (825.0,45.0) ellipse (5.0 and 5.0 );
\fill [line width=3.0pt, custcol_0_0_0]  (975.0,5.0) ellipse (5.0 and 5.0 );
\fill [line width=3.0pt, custcol_0_0_0]  (875.0,45.0) ellipse (5.0 and 5.0 );
\fill [line width=3.0pt, custcol_0_0_0]  (1025.0,5.0) ellipse (5.0 and 5.0 );
\fill [line width=3.0pt, custcol_0_0_0]  (545.0,45.0) ellipse (5.0 and 5.0 );
\fill [line width=3.0pt, custcol_0_0_0]  (265.0,25.0) ellipse (5.0 and 5.0 );
\fill [line width=3.0pt, custcol_0_0_0]  (1075.0,5.0) ellipse (5.0 and 5.0 );
\fill [line width=3.0pt, custcol_0_0_0]  (645.0,25.0) ellipse (5.0 and 5.0 );
\fill [line width=3.0pt, custcol_0_0_0]  (365.0,5.0) ellipse (5.0 and 5.0 );
\fill [line width=3.0pt, custcol_0_0_0]  (695.0,25.0) ellipse (5.0 and 5.0 );
\fill [line width=3.0pt, custcol_0_0_0]  (745.0,25.0) ellipse (5.0 and 5.0 );
\fill [line width=3.0pt, custcol_0_0_0]  (845.0,5.0) ellipse (5.0 and 5.0 );
\fill [line width=3.0pt, custcol_0_0_0]  (415.0,45.0) ellipse (5.0 and 5.0 );
\fill [line width=3.0pt, custcol_0_0_0]  (565.0,5.0) ellipse (5.0 and 5.0 );
\fill [line width=3.0pt, custcol_0_0_0]  (895.0,25.0) ellipse (5.0 and 5.0 );
\fill [line width=3.0pt, custcol_0_0_0]  (465.0,45.0) ellipse (5.0 and 5.0 );
\fill [line width=3.0pt, custcol_0_0_0]  (795.0,65.0) ellipse (5.0 and 5.0 );
\fill [line width=3.0pt, custcol_0_0_0]  (285.0,5.0) ellipse (5.0 and 5.0 );
\fill [line width=3.0pt, custcol_0_0_0]  (615.0,45.0) ellipse (5.0 and 5.0 );
\fill [line width=3.0pt, custcol_0_0_0]  (435.0,5.0) ellipse (5.0 and 5.0 );
\fill [line width=3.0pt, custcol_0_0_0]  (665.0,45.0) ellipse (5.0 and 5.0 );
\fill [line width=3.0pt, custcol_0_0_0]  (335.0,45.0) ellipse (5.0 and 5.0 );
\fill [line width=3.0pt, custcol_0_0_0]  (485.0,5.0) ellipse (5.0 and 5.0 );
\fill [line width=3.0pt, custcol_0_0_0]  (715.0,45.0) ellipse (5.0 and 5.0 );
\fill [line width=3.0pt, custcol_0_0_0]  (815.0,25.0) ellipse (5.0 and 5.0 );
\fill [line width=3.0pt, custcol_0_0_0]  (385.0,45.0) ellipse (5.0 and 5.0 );
\fill [line width=3.0pt, custcol_0_0_0]  (1095.0,65.0) ellipse (5.0 and 5.0 );
\fill [line width=3.0pt, custcol_0_0_0]  (535.0,25.0) ellipse (5.0 and 5.0 );
\fill [line width=3.0pt, custcol_0_0_0]  (1145.0,65.0) ellipse (5.0 and 5.0 );
\fill [line width=3.0pt, custcol_0_0_0]  (635.0,5.0) ellipse (5.0 and 5.0 );
\fill [line width=3.0pt, custcol_0_0_0]  (1195.0,65.0) ellipse (5.0 and 5.0 );
\fill [line width=3.0pt, custcol_0_0_0]  (685.0,5.0) ellipse (5.0 and 5.0 );
\fill [line width=3.0pt, custcol_0_0_0]  (1245.0,65.0) ellipse (5.0 and 5.0 );
\fill [line width=3.0pt, custcol_0_0_0]  (735.0,5.0) ellipse (5.0 and 5.0 );
\fill [line width=3.0pt, custcol_0_0_0]  (915.0,65.0) ellipse (5.0 and 5.0 );
\fill [line width=3.0pt, custcol_0_0_0]  (1015.0,45.0) ellipse (5.0 and 5.0 );
\fill [line width=3.0pt, custcol_0_0_0]  (585.0,65.0) ellipse (5.0 and 5.0 );
\fill [line width=3.0pt, custcol_0_0_0]  (305.0,45.0) ellipse (5.0 and 5.0 );
\fill [line width=3.0pt, custcol_0_0_0]  (965.0,65.0) ellipse (5.0 and 5.0 );
\fill [line width=3.0pt, custcol_0_0_0]  (1115.0,25.0) ellipse (5.0 and 5.0 );
\fill [line width=3.0pt, custcol_0_0_0]  (405.0,25.0) ellipse (5.0 and 5.0 );
\fill [line width=3.0pt, custcol_0_0_0]  (1165.0,25.0) ellipse (5.0 and 5.0 );
\fill [line width=3.0pt, custcol_0_0_0]  (1065.0,65.0) ellipse (5.0 and 5.0 );
\fill [line width=3.0pt, custcol_0_0_0]  (1215.0,25.0) ellipse (5.0 and 5.0 );
\fill [line width=3.0pt, custcol_0_0_0]  (785.0,45.0) ellipse (5.0 and 5.0 );
\fill [line width=3.0pt, custcol_0_0_0]  (935.0,25.0) ellipse (5.0 and 5.0 );
\fill [line width=3.0pt, custcol_0_0_0]  (505.0,45.0) ellipse (5.0 and 5.0 );
\fill [line width=3.0pt, custcol_0_0_0]  (605.0,25.0) ellipse (5.0 and 5.0 );
\fill [line width=3.0pt, custcol_0_0_0]  (835.0,65.0) ellipse (5.0 and 5.0 );
\fill [line width=3.0pt, custcol_0_0_0]  (325.0,5.0) ellipse (5.0 and 5.0 );
\fill [line width=3.0pt, custcol_0_0_0]  (985.0,25.0) ellipse (5.0 and 5.0 );
\fill [line width=3.0pt, custcol_0_0_0]  (885.0,65.0) ellipse (5.0 and 5.0 );
\fill [line width=3.0pt, custcol_0_0_0]  (1035.0,25.0) ellipse (5.0 and 5.0 );
\fill [line width=3.0pt, custcol_0_0_0]  (375.0,25.0) ellipse (5.0 and 5.0 );
\fill [line width=3.0pt, custcol_0_0_0]  (525.0,5.0) ellipse (5.0 and 5.0 );
\fill [line width=3.0pt, custcol_0_0_0]  (1135.0,45.0) ellipse (5.0 and 5.0 );
\fill [line width=3.0pt, custcol_0_0_0]  (1185.0,45.0) ellipse (5.0 and 5.0 );
\fill [line width=3.0pt, custcol_0_0_0]  (755.0,65.0) ellipse (5.0 and 5.0 );
\fill [line width=3.0pt, custcol_0_0_0]  (245.0,5.0) ellipse (5.0 and 5.0 );
\fill [line width=3.0pt, custcol_0_0_0]  (855.0,45.0) ellipse (5.0 and 5.0 );
\fill [line width=3.0pt, custcol_0_0_0]  (425.0,65.0) ellipse (5.0 and 5.0 );
\fill [line width=3.0pt, custcol_0_0_0]  (1235.0,45.0) ellipse (5.0 and 5.0 );
\fill [line width=3.0pt, custcol_0_0_0]  (905.0,45.0) ellipse (5.0 and 5.0 );
\fill [line width=3.0pt, custcol_0_0_0]  (475.0,65.0) ellipse (5.0 and 5.0 );
\fill [line width=3.0pt, custcol_0_0_0]  (295.0,25.0) ellipse (5.0 and 5.0 );
\fill [line width=3.0pt, custcol_0_0_0]  (955.0,45.0) ellipse (5.0 and 5.0 );
\fill [line width=3.0pt, custcol_0_0_0]  (1105.0,5.0) ellipse (5.0 and 5.0 );
\fill [line width=3.0pt, custcol_0_0_0]  (1005.0,45.0) ellipse (5.0 and 5.0 );
\fill [line width=3.0pt, custcol_0_0_0]  (1155.0,5.0) ellipse (5.0 and 5.0 );
\draw [line width=3.0pt, custcol_0_0_0 ]  (245.0,95.0) -- (245.0,5.0) -- (235.0,5.0) -- (221.0,5.0) ;
\draw [line width=3.0pt, custcol_0_0_0 ]  (265.0,25.0) -- (235.0,25.0) -- (221.0,25.0) ;
\draw [line width=3.0pt, custcol_0_0_0 ]  (305.0,95.0) -- (305.0,45.0) -- (235.0,45.0) -- (221.0,45.0) ;
\draw [line width=3.0pt, custcol_0_0_0 ]  (475.0,65.0) -- (425.0,65.0) -- (345.0,65.0) -- (235.0,65.0) -- (221.0,65.0) ;
\draw [line width=5.0pt, custcol_0_0_0 ]  (216.0,45.0) -- (221.0,45.0) ;
\draw [line width=5.0pt, custcol_0_0_0 ]  (221.0,6.0) -- (221.0,64.0) ;
\logisimfontA{\fontsize{7pt}{7pt}\selectfont\node[inner sep=0, outer sep=0, custcol_0_0_0, anchor=base west] at  (224.0,2.0)  {3};}
\logisimfontA{\fontsize{7pt}{7pt}\selectfont\node[inner sep=0, outer sep=0, custcol_0_0_0, anchor=base west] at  (224.0,22.0)  {2};}
\logisimfontA{\fontsize{7pt}{7pt}\selectfont\node[inner sep=0, outer sep=0, custcol_0_0_0, anchor=base west] at  (224.0,42.0)  {1};}
\logisimfontA{\fontsize{7pt}{7pt}\selectfont\node[inner sep=0, outer sep=0, custcol_0_0_0, anchor=base west] at  (224.0,62.0)  {0};}
\fill [line width=5.0pt, custcol_0_0_0]  (215.0,45.0) ellipse (2.0 and 2.0 );
\fill [line width=5.0pt, custcol_0_0_0]  (235.0,5.0) ellipse (2.0 and 2.0 );
\fill [line width=5.0pt, custcol_0_0_0]  (235.0,25.0) ellipse (2.0 and 2.0 );
\fill [line width=5.0pt, custcol_0_0_0]  (235.0,45.0) ellipse (2.0 and 2.0 );
\fill [line width=5.0pt, custcol_0_0_0]  (235.0,65.0) ellipse (2.0 and 2.0 );
\draw [line width=4.0pt, custcol_0_0_0 ]  (182.0,45.0) -- (185.0,45.0) -- (215.0,45.0) ;
\draw [line width=2.0pt, custcol_0_0_0 ]  (170.0,54.0) -- (180.0,45.0) -- (170.0,36.0) -- (116.0,36.0) -- (116.0,54.0) -- cycle;
\logisimfontB{\fontsize{12pt}{12pt}\selectfont\node[inner sep=0, outer sep=0, custcol_0_0_0, anchor=base west] at  (136.0,52.0)  {x4};}
\logisimfontC{\fontsize{16pt}{16pt}\fontseries{bx}\selectfont\node[inner sep=0, outer sep=0, custcol_0_0_0, anchor=base west] at  (100.0,25.0)  {CurrentState};}
\fill [line width=2.0pt, custcol_0_0_0]  (185.0,45.0) ellipse (2.0 and 2.0 );
\draw [line width=2.0pt, custcol_0_0_0]  (325.0,100.0) ellipse (4.5 and 4.5 );
\draw [line width=2.0pt, custcol_0_0_0]  (345.0,100.0) ellipse (4.5 and 4.5 );
\draw [line width=2.0pt, custcol_0_0_0] (320.0,120.0) arc (-180.0:-360.0:15.0 and 15.0 );
\draw [line width=2.0pt, custcol_0_0_0 ]  (350.0,120.0) -- (350.0,106.0) -- (320.0,106.0) -- (320.0,120.0) ;
\draw [line width=2.0pt, custcol_0_0_0]  (375.0,100.0) ellipse (4.5 and 4.5 );
\draw [line width=2.0pt, custcol_0_0_0]  (385.0,100.0) ellipse (4.5 and 4.5 );
\draw [line width=2.0pt, custcol_0_0_0] (360.0,120.0) arc (-180.0:-360.0:15.0 and 15.0 );
\draw [line width=2.0pt, custcol_0_0_0 ]  (390.0,120.0) -- (390.0,106.0) -- (360.0,106.0) -- (360.0,120.0) ;
\draw [line width=3.0pt, custcol_0_0_0 ]  (255.0,135.0) -- (255.0,145.0) -- (265.0,145.0) -- (265.0,155.0) -- (265.0,157.0) ;
\draw [line width=3.0pt, custcol_0_0_0 ]  (295.0,135.0) -- (295.0,145.0) -- (285.0,145.0) -- (285.0,155.0) -- (285.0,157.0) ;
\draw [line width=2.0pt, custcol_0_0_0 ]  (275.0,185.0) .. controls  (290.0,175.0)  ..  (290.0,155.0) .. controls  (275.0,163.0)  ..  (260.0,155.0) .. controls  (260.0,175.0)  ..  (275.0,185.0) -- cycle ;
\draw [line width=2.0pt, custcol_0_0_0]  (685.0,100.0) ellipse (4.5 and 4.5 );
\draw [line width=2.0pt, custcol_0_0_0]  (695.0,100.0) ellipse (4.5 and 4.5 );
\draw [line width=2.0pt, custcol_0_0_0]  (715.0,100.0) ellipse (4.5 and 4.5 );
\draw [line width=2.0pt, custcol_0_0_0]  (725.0,100.0) ellipse (4.5 and 4.5 );
\draw [line width=2.0pt, custcol_0_0_0] (690.0,120.0) arc (-180.0:-360.0:15.0 and 15.0 );
\draw [line width=2.0pt, custcol_0_0_0 ]  (720.0,120.0) -- (720.0,106.0) -- (690.0,106.0) -- (690.0,120.0) ;
\draw [line width=2.0pt, custcol_0_0_0]  (745.0,100.0) ellipse (4.5 and 4.5 );
\draw [line width=2.0pt, custcol_0_0_0] (730.0,120.0) arc (-180.0:-360.0:15.0 and 15.0 );
\draw [line width=2.0pt, custcol_0_0_0 ]  (760.0,120.0) -- (760.0,106.0) -- (730.0,106.0) -- (730.0,120.0) ;
\draw [line width=3.0pt, custcol_0_0_0 ]  (615.0,135.0) -- (615.0,155.0) -- (615.0,157.0) ;
\draw [line width=3.0pt, custcol_0_0_0 ]  (655.0,135.0) -- (655.0,145.0) -- (635.0,145.0) -- (635.0,155.0) -- (635.0,157.0) ;
\draw [line width=2.0pt, custcol_0_0_0 ]  (625.0,185.0) .. controls  (640.0,175.0)  ..  (640.0,155.0) .. controls  (625.0,163.0)  ..  (610.0,155.0) .. controls  (610.0,175.0)  ..  (625.0,185.0) -- cycle ;
\draw [line width=2.0pt, custcol_0_0_0]  (245.0,100.0) ellipse (4.5 and 4.5 );
\draw [line width=2.0pt, custcol_0_0_0] (240.0,120.0) arc (-180.0:-360.0:15.0 and 15.0 );
\draw [line width=2.0pt, custcol_0_0_0 ]  (270.0,120.0) -- (270.0,106.0) -- (240.0,106.0) -- (240.0,120.0) ;
\draw [line width=2.0pt, custcol_0_0_0]  (535.0,100.0) ellipse (4.5 and 4.5 );
\draw [line width=2.0pt, custcol_0_0_0]  (545.0,100.0) ellipse (4.5 and 4.5 );
\draw [line width=2.0pt, custcol_0_0_0] (520.0,120.0) arc (-180.0:-360.0:15.0 and 15.0 );
\draw [line width=2.0pt, custcol_0_0_0 ]  (550.0,120.0) -- (550.0,106.0) -- (520.0,106.0) -- (520.0,120.0) ;
\draw [line width=2.0pt, custcol_0_0_0]  (435.0,100.0) ellipse (4.5 and 4.5 );
\draw [line width=2.0pt, custcol_0_0_0]  (445.0,100.0) ellipse (4.5 and 4.5 );
\draw [line width=2.0pt, custcol_0_0_0]  (465.0,100.0) ellipse (4.5 and 4.5 );
\draw [line width=2.0pt, custcol_0_0_0]  (475.0,100.0) ellipse (4.5 and 4.5 );
\draw [line width=2.0pt, custcol_0_0_0] (440.0,120.0) arc (-180.0:-360.0:15.0 and 15.0 );
\draw [line width=2.0pt, custcol_0_0_0 ]  (470.0,120.0) -- (470.0,106.0) -- (440.0,106.0) -- (440.0,120.0) ;
\draw [line width=2.0pt, custcol_0_0_0 ]  (480.0,106.0) -- (430.0,106.0) ;
\draw [line width=2.0pt, custcol_0_0_0]  (295.0,100.0) ellipse (4.5 and 4.5 );
\draw [line width=2.0pt, custcol_0_0_0]  (305.0,100.0) ellipse (4.5 and 4.5 );
\draw [line width=2.0pt, custcol_0_0_0] (280.0,120.0) arc (-180.0:-360.0:15.0 and 15.0 );
\draw [line width=2.0pt, custcol_0_0_0 ]  (310.0,120.0) -- (310.0,106.0) -- (280.0,106.0) -- (280.0,120.0) ;
\draw [line width=2.0pt, custcol_0_0_0]  (565.0,100.0) ellipse (4.5 and 4.5 );
\draw [line width=2.0pt, custcol_0_0_0] (560.0,120.0) arc (-180.0:-360.0:15.0 and 15.0 );
\draw [line width=2.0pt, custcol_0_0_0 ]  (590.0,120.0) -- (590.0,106.0) -- (560.0,106.0) -- (560.0,120.0) ;
\draw [line width=3.0pt, custcol_0_0_0 ]  (495.0,135.0) -- (495.0,145.0) -- (485.0,145.0) -- (485.0,155.0) -- (485.0,157.0) ;
\draw [line width=3.0pt, custcol_0_0_0 ]  (535.0,135.0) -- (535.0,145.0) -- (505.0,145.0) -- (505.0,155.0) -- (505.0,157.0) ;
\draw [line width=2.0pt, custcol_0_0_0 ]  (495.0,185.0) .. controls  (510.0,175.0)  ..  (510.0,155.0) .. controls  (495.0,163.0)  ..  (480.0,155.0) .. controls  (480.0,175.0)  ..  (495.0,185.0) -- cycle ;
\draw [line width=2.0pt, custcol_0_0_0]  (485.0,100.0) ellipse (4.5 and 4.5 );
\draw [line width=2.0pt, custcol_0_0_0] (480.0,120.0) arc (-180.0:-360.0:15.0 and 15.0 );
\draw [line width=2.0pt, custcol_0_0_0 ]  (510.0,120.0) -- (510.0,106.0) -- (480.0,106.0) -- (480.0,120.0) ;
\draw [line width=2.0pt, custcol_0_0_0]  (615.0,100.0) ellipse (4.5 and 4.5 );
\draw [line width=2.0pt, custcol_0_0_0] (600.0,120.0) arc (-180.0:-360.0:15.0 and 15.0 );
\draw [line width=2.0pt, custcol_0_0_0 ]  (630.0,120.0) -- (630.0,106.0) -- (600.0,106.0) -- (600.0,120.0) ;
\draw [line width=3.0pt, custcol_0_0_0 ]  (415.0,135.0) -- (415.0,145.0) -- (445.0,145.0) -- (445.0,155.0) -- (445.0,157.0) ;
\draw [line width=3.0pt, custcol_0_0_0 ]  (455.0,135.0) -- (455.0,145.0) -- (465.0,145.0) -- (465.0,155.0) -- (465.0,157.0) ;
\draw [line width=2.0pt, custcol_0_0_0 ]  (455.0,185.0) .. controls  (470.0,175.0)  ..  (470.0,155.0) .. controls  (455.0,163.0)  ..  (440.0,155.0) .. controls  (440.0,175.0)  ..  (455.0,185.0) -- cycle ;
\draw [line width=2.0pt, custcol_0_0_0]  (415.0,100.0) ellipse (4.5 and 4.5 );
\draw [line width=2.0pt, custcol_0_0_0] (400.0,120.0) arc (-180.0:-360.0:15.0 and 15.0 );
\draw [line width=2.0pt, custcol_0_0_0 ]  (430.0,120.0) -- (430.0,106.0) -- (400.0,106.0) -- (400.0,120.0) ;
\draw [line width=2.0pt, custcol_0_0_0]  (635.0,100.0) ellipse (4.5 and 4.5 );
\draw [line width=2.0pt, custcol_0_0_0]  (645.0,100.0) ellipse (4.5 and 4.5 );
\draw [line width=2.0pt, custcol_0_0_0]  (665.0,100.0) ellipse (4.5 and 4.5 );
\draw [line width=2.0pt, custcol_0_0_0]  (675.0,100.0) ellipse (4.5 and 4.5 );
\draw [line width=2.0pt, custcol_0_0_0] (640.0,120.0) arc (-180.0:-360.0:15.0 and 15.0 );
\draw [line width=2.0pt, custcol_0_0_0 ]  (670.0,120.0) -- (670.0,106.0) -- (640.0,106.0) -- (640.0,120.0) ;
\draw [line width=2.0pt, custcol_0_0_0 ]  (730.0,106.0) -- (680.0,106.0) -- (630.0,106.0) ;
\draw [line width=2.0pt, custcol_0_0_0]  (785.0,100.0) ellipse (4.5 and 4.5 );
\draw [line width=2.0pt, custcol_0_0_0] (770.0,120.0) arc (-180.0:-360.0:15.0 and 15.0 );
\draw [line width=2.0pt, custcol_0_0_0 ]  (800.0,120.0) -- (800.0,106.0) -- (770.0,106.0) -- (770.0,120.0) ;
\draw [line width=3.0pt, custcol_0_0_0 ]  (745.0,135.0) -- (745.0,145.0) -- (755.0,145.0) -- (755.0,155.0) -- (755.0,157.0) ;
\draw [line width=3.0pt, custcol_0_0_0 ]  (785.0,135.0) -- (785.0,145.0) -- (775.0,145.0) -- (775.0,155.0) -- (775.0,157.0) ;
\draw [line width=2.0pt, custcol_0_0_0 ]  (765.0,185.0) .. controls  (780.0,175.0)  ..  (780.0,155.0) .. controls  (765.0,163.0)  ..  (750.0,155.0) .. controls  (750.0,175.0)  ..  (765.0,185.0) -- cycle ;
\draw [line width=3.0pt, custcol_0_0_0 ]  (765.0,239.0) -- (765.0,236.0) ;
\draw [line width=2.0pt, custcol_0_0_0 ]  (774.0,265.0) -- (765.0,275.0) -- (756.0,265.0) -- (756.0,241.0) -- (774.0,241.0) -- cycle;
\logisimfontB{\fontsize{12pt}{12pt}\selectfont\node[inner sep=0, outer sep=0, custcol_0_0_0, anchor=base west, rotate=90.0] at  (772.0,270.0)  {x1};}
\logisimfontC{\fontsize{10pt}{10pt}\fontseries{bx}\selectfont\node[inner sep=0, outer sep=0, custcol_0_0_0, anchor=base west] at  (740.0,287.0)  {RegWrite};}
\fill [line width=2.0pt, custcol_0_0_0]  (765.0,235.0) ellipse (2.0 and 2.0 );
\draw [line width=2.0pt, custcol_0_0_0]  (815.0,100.0) ellipse (4.5 and 4.5 );
\draw [line width=2.0pt, custcol_0_0_0] (810.0,120.0) arc (-180.0:-360.0:15.0 and 15.0 );
\draw [line width=2.0pt, custcol_0_0_0 ]  (840.0,120.0) -- (840.0,106.0) -- (810.0,106.0) -- (810.0,120.0) ;
\draw [line width=2.0pt, custcol_0_0_0]  (875.0,100.0) ellipse (4.5 and 4.5 );
\draw [line width=2.0pt, custcol_0_0_0] (850.0,120.0) arc (-180.0:-360.0:15.0 and 15.0 );
\draw [line width=2.0pt, custcol_0_0_0 ]  (880.0,120.0) -- (880.0,106.0) -- (850.0,106.0) -- (850.0,120.0) ;
\draw [line width=2.0pt, custcol_0_0_0 ]  (890.0,106.0) -- (840.0,106.0) ;
\draw [line width=3.0pt, custcol_0_0_0 ]  (825.0,135.0) -- (825.0,145.0) -- (835.0,145.0) -- (835.0,155.0) -- (835.0,157.0) ;
\draw [line width=3.0pt, custcol_0_0_0 ]  (865.0,135.0) -- (865.0,145.0) -- (855.0,145.0) -- (855.0,155.0) -- (855.0,157.0) ;
\draw [line width=2.0pt, custcol_0_0_0 ]  (845.0,185.0) .. controls  (860.0,175.0)  ..  (860.0,155.0) .. controls  (845.0,163.0)  ..  (830.0,155.0) .. controls  (830.0,175.0)  ..  (845.0,185.0) -- cycle ;
\draw [line width=3.0pt, custcol_0_0_0 ]  (905.0,135.0) -- (905.0,165.0) -- (885.0,165.0) -- (885.0,195.0) -- (885.0,209.0) ;
\draw [line width=3.0pt, custcol_0_0_0 ]  (845.0,185.0) -- (845.0,195.0) -- (845.0,209.0) ;
\draw [line width=5.0pt, custcol_0_0_0 ]  (885.0,214.0) -- (885.0,209.0) ;
\draw [line width=5.0pt, custcol_0_0_0 ]  (884.0,209.0) -- (846.0,209.0) ;
\logisimfontA{\fontsize{7pt}{7pt}\selectfont\node[inner sep=0, outer sep=0, custcol_0_0_0, anchor=base west, rotate=-90.0] at  (888.0,202.0)  {0};}
\logisimfontA{\fontsize{7pt}{7pt}\selectfont\node[inner sep=0, outer sep=0, custcol_0_0_0, anchor=base west, rotate=-90.0] at  (848.0,202.0)  {1};}
\fill [line width=5.0pt, custcol_0_0_0]  (885.0,215.0) ellipse (2.0 and 2.0 );
\fill [line width=5.0pt, custcol_0_0_0]  (885.0,195.0) ellipse (2.0 and 2.0 );
\fill [line width=5.0pt, custcol_0_0_0]  (845.0,195.0) ellipse (2.0 and 2.0 );
\draw [line width=2.0pt, custcol_0_0_0]  (935.0,100.0) ellipse (4.5 and 4.5 );
\draw [line width=2.0pt, custcol_0_0_0]  (955.0,100.0) ellipse (4.5 and 4.5 );
\draw [line width=2.0pt, custcol_0_0_0]  (965.0,100.0) ellipse (4.5 and 4.5 );
\draw [line width=2.0pt, custcol_0_0_0] (930.0,120.0) arc (-180.0:-360.0:15.0 and 15.0 );
\draw [line width=2.0pt, custcol_0_0_0 ]  (960.0,120.0) -- (960.0,106.0) -- (930.0,106.0) -- (930.0,120.0) ;
\draw [line width=2.0pt, custcol_0_0_0]  (905.0,100.0) ellipse (4.5 and 4.5 );
\draw [line width=2.0pt, custcol_0_0_0] (890.0,120.0) arc (-180.0:-360.0:15.0 and 15.0 );
\draw [line width=2.0pt, custcol_0_0_0 ]  (920.0,120.0) -- (920.0,106.0) -- (890.0,106.0) -- (890.0,120.0) ;
\draw [line width=2.0pt, custcol_0_0_0]  (975.0,100.0) ellipse (4.5 and 4.5 );
\draw [line width=2.0pt, custcol_0_0_0]  (985.0,100.0) ellipse (4.5 and 4.5 );
\draw [line width=2.0pt, custcol_0_0_0]  (1005.0,100.0) ellipse (4.5 and 4.5 );
\draw [line width=2.0pt, custcol_0_0_0]  (1015.0,100.0) ellipse (4.5 and 4.5 );
\draw [line width=2.0pt, custcol_0_0_0] (980.0,120.0) arc (-180.0:-360.0:15.0 and 15.0 );
\draw [line width=2.0pt, custcol_0_0_0 ]  (1010.0,120.0) -- (1010.0,106.0) -- (980.0,106.0) -- (980.0,120.0) ;
\draw [line width=2.0pt, custcol_0_0_0]  (1025.0,100.0) ellipse (4.5 and 4.5 );
\draw [line width=2.0pt, custcol_0_0_0]  (1055.0,100.0) ellipse (4.5 and 4.5 );
\draw [line width=2.0pt, custcol_0_0_0] (1030.0,120.0) arc (-180.0:-360.0:15.0 and 15.0 );
\draw [line width=2.0pt, custcol_0_0_0 ]  (1060.0,120.0) -- (1060.0,106.0) -- (1030.0,106.0) -- (1030.0,120.0) ;
\draw [line width=2.0pt, custcol_0_0_0 ]  (1070.0,106.0) -- (1020.0,106.0) -- (970.0,106.0) -- (920.0,106.0) ;
\draw [line width=3.0pt, custcol_0_0_0 ]  (995.0,135.0) -- (995.0,145.0) -- (1005.0,145.0) -- (1005.0,155.0) -- (1005.0,157.0) ;
\draw [line width=3.0pt, custcol_0_0_0 ]  (1045.0,135.0) -- (1045.0,145.0) -- (1025.0,145.0) -- (1025.0,155.0) -- (1025.0,157.0) ;
\draw [line width=2.0pt, custcol_0_0_0 ]  (1015.0,185.0) .. controls  (1030.0,175.0)  ..  (1030.0,155.0) .. controls  (1015.0,163.0)  ..  (1000.0,155.0) .. controls  (1000.0,175.0)  ..  (1015.0,185.0) -- cycle ;
\draw [line width=4.0pt, custcol_0_0_0 ]  (865.0,229.0) -- (865.0,228.0) ;
\draw [line width=2.0pt, custcol_0_0_0 ]  (874.0,265.0) -- (865.0,275.0) -- (856.0,265.0) -- (856.0,231.0) -- (874.0,231.0) -- cycle;
\logisimfontB{\fontsize{12pt}{12pt}\selectfont\node[inner sep=0, outer sep=0, custcol_0_0_0, anchor=base west, rotate=90.0] at  (872.0,265.0)  {x2};}
\logisimfontC{\fontsize{10pt}{10pt}\fontseries{bx}\selectfont\node[inner sep=0, outer sep=0, custcol_0_0_0, anchor=base west] at  (839.0,287.0)  {Mem2Reg};}
\fill [line width=2.0pt, custcol_0_0_0]  (865.0,225.0) ellipse (2.0 and 2.0 );
\draw [line width=2.0pt, custcol_0_0_0]  (1105.0,100.0) ellipse (4.5 and 4.5 );
\draw [line width=2.0pt, custcol_0_0_0] (1110.0,120.0) arc (-180.0:-360.0:15.0 and 15.0 );
\draw [line width=2.0pt, custcol_0_0_0 ]  (1140.0,120.0) -- (1140.0,106.0) -- (1110.0,106.0) -- (1110.0,120.0) ;
\draw [line width=2.0pt, custcol_0_0_0]  (1075.0,100.0) ellipse (4.5 and 4.5 );
\draw [line width=2.0pt, custcol_0_0_0] (1070.0,120.0) arc (-180.0:-360.0:15.0 and 15.0 );
\draw [line width=2.0pt, custcol_0_0_0 ]  (1100.0,120.0) -- (1100.0,106.0) -- (1070.0,106.0) -- (1070.0,120.0) ;
\draw [line width=3.0pt, custcol_0_0_0 ]  (1125.0,239.0) -- (1125.0,236.0) ;
\draw [line width=2.0pt, custcol_0_0_0 ]  (1134.0,265.0) -- (1125.0,275.0) -- (1116.0,265.0) -- (1116.0,241.0) -- (1134.0,241.0) -- cycle;
\logisimfontB{\fontsize{12pt}{12pt}\selectfont\node[inner sep=0, outer sep=0, custcol_0_0_0, anchor=base west, rotate=90.0] at  (1132.0,270.0)  {x1};}
\logisimfontC{\fontsize{10pt}{10pt}\fontseries{bx}\selectfont\node[inner sep=0, outer sep=0, custcol_0_0_0, anchor=base west] at  (1098.0,287.0)  {MemWrite};}
\fill [line width=2.0pt, custcol_0_0_0]  (1125.0,235.0) ellipse (2.0 and 2.0 );
\draw [line width=2.0pt, custcol_0_0_0]  (1155.0,100.0) ellipse (4.5 and 4.5 );
\draw [line width=2.0pt, custcol_0_0_0]  (1165.0,100.0) ellipse (4.5 and 4.5 );
\draw [line width=2.0pt, custcol_0_0_0]  (1185.0,100.0) ellipse (4.5 and 4.5 );
\draw [line width=2.0pt, custcol_0_0_0]  (1195.0,100.0) ellipse (4.5 and 4.5 );
\draw [line width=2.0pt, custcol_0_0_0] (1160.0,120.0) arc (-180.0:-360.0:15.0 and 15.0 );
\draw [line width=2.0pt, custcol_0_0_0 ]  (1190.0,120.0) -- (1190.0,106.0) -- (1160.0,106.0) -- (1160.0,120.0) ;
\draw [line width=2.0pt, custcol_0_0_0]  (1205.0,100.0) ellipse (4.5 and 4.5 );
\draw [line width=2.0pt, custcol_0_0_0]  (1245.0,100.0) ellipse (4.5 and 4.5 );
\draw [line width=2.0pt, custcol_0_0_0] (1210.0,120.0) arc (-180.0:-360.0:15.0 and 15.0 );
\draw [line width=2.0pt, custcol_0_0_0 ]  (1240.0,120.0) -- (1240.0,106.0) -- (1210.0,106.0) -- (1210.0,120.0) ;
\draw [line width=3.0pt, custcol_0_0_0 ]  (1175.0,135.0) -- (1175.0,145.0) -- (1195.0,145.0) -- (1195.0,155.0) -- (1195.0,157.0) ;
\draw [line width=3.0pt, custcol_0_0_0 ]  (1225.0,135.0) -- (1225.0,145.0) -- (1215.0,145.0) -- (1215.0,155.0) -- (1215.0,157.0) ;
\draw [line width=2.0pt, custcol_0_0_0 ]  (1205.0,185.0) .. controls  (1220.0,175.0)  ..  (1220.0,155.0) .. controls  (1205.0,163.0)  ..  (1190.0,155.0) .. controls  (1190.0,175.0)  ..  (1205.0,185.0) -- cycle ;
\draw [line width=2.0pt, custcol_0_0_0]  (1255.0,100.0) ellipse (4.5 and 4.5 );
\draw [line width=2.0pt, custcol_0_0_0]  (1265.0,100.0) ellipse (4.5 and 4.5 );
\draw [line width=2.0pt, custcol_0_0_0]  (1285.0,100.0) ellipse (4.5 and 4.5 );
\draw [line width=2.0pt, custcol_0_0_0]  (1295.0,100.0) ellipse (4.5 and 4.5 );
\draw [line width=2.0pt, custcol_0_0_0] (1260.0,120.0) arc (-180.0:-360.0:15.0 and 15.0 );
\draw [line width=2.0pt, custcol_0_0_0 ]  (1290.0,120.0) -- (1290.0,106.0) -- (1260.0,106.0) -- (1260.0,120.0) ;
\draw [line width=2.0pt, custcol_0_0_0 ]  (1300.0,106.0) -- (1250.0,106.0) -- (1200.0,106.0) -- (1150.0,106.0) -- (1100.0,106.0) ;
\draw [line width=3.0pt, custcol_0_0_0 ]  (945.0,239.0) -- (945.0,236.0) ;
\draw [line width=2.0pt, custcol_0_0_0 ]  (954.0,265.0) -- (945.0,275.0) -- (936.0,265.0) -- (936.0,241.0) -- (954.0,241.0) -- cycle;
\logisimfontB{\fontsize{12pt}{12pt}\selectfont\node[inner sep=0, outer sep=0, custcol_0_0_0, anchor=base west, rotate=90.0] at  (952.0,270.0)  {x1};}
\logisimfontC{\fontsize{10pt}{10pt}\fontseries{bx}\selectfont\node[inner sep=0, outer sep=0, custcol_0_0_0, anchor=base west] at  (908.0,287.0)  {PCWriteCond};}
\fill [line width=2.0pt, custcol_0_0_0]  (945.0,235.0) ellipse (2.0 and 2.0 );
\draw [line width=3.0pt, custcol_0_0_0 ]  (1015.0,239.0) -- (1015.0,236.0) ;
\draw [line width=2.0pt, custcol_0_0_0 ]  (1024.0,265.0) -- (1015.0,275.0) -- (1006.0,265.0) -- (1006.0,241.0) -- (1024.0,241.0) -- cycle;
\logisimfontB{\fontsize{12pt}{12pt}\selectfont\node[inner sep=0, outer sep=0, custcol_0_0_0, anchor=base west, rotate=90.0] at  (1022.0,270.0)  {x1};}
\logisimfontC{\fontsize{10pt}{10pt}\fontseries{bx}\selectfont\node[inner sep=0, outer sep=0, custcol_0_0_0, anchor=base west] at  (993.0,287.0)  {PCWrite};}
\fill [line width=2.0pt, custcol_0_0_0]  (1015.0,235.0) ellipse (2.0 and 2.0 );
\draw [line width=3.0pt, custcol_0_0_0 ]  (1065.0,239.0) -- (1065.0,236.0) ;
\draw [line width=2.0pt, custcol_0_0_0 ]  (1074.0,265.0) -- (1065.0,275.0) -- (1056.0,265.0) -- (1056.0,241.0) -- (1074.0,241.0) -- cycle;
\logisimfontB{\fontsize{12pt}{12pt}\selectfont\node[inner sep=0, outer sep=0, custcol_0_0_0, anchor=base west, rotate=90.0] at  (1072.0,270.0)  {x1};}
\logisimfontC{\fontsize{10pt}{10pt}\fontseries{bx}\selectfont\node[inner sep=0, outer sep=0, custcol_0_0_0, anchor=base west] at  (1053.0,287.0)  {IorD};}
\fill [line width=2.0pt, custcol_0_0_0]  (1065.0,235.0) ellipse (2.0 and 2.0 );
\draw [line width=3.0pt, custcol_0_0_0 ]  (1205.0,239.0) -- (1205.0,236.0) ;
\draw [line width=2.0pt, custcol_0_0_0 ]  (1214.0,265.0) -- (1205.0,275.0) -- (1196.0,265.0) -- (1196.0,241.0) -- (1214.0,241.0) -- cycle;
\logisimfontB{\fontsize{12pt}{12pt}\selectfont\node[inner sep=0, outer sep=0, custcol_0_0_0, anchor=base west, rotate=90.0] at  (1212.0,270.0)  {x1};}
\logisimfontC{\fontsize{10pt}{10pt}\fontseries{bx}\selectfont\node[inner sep=0, outer sep=0, custcol_0_0_0, anchor=base west] at  (1179.0,287.0)  {MemRead};}
\fill [line width=2.0pt, custcol_0_0_0]  (1205.0,235.0) ellipse (2.0 and 2.0 );
\draw [line width=3.0pt, custcol_0_0_0 ]  (1275.0,239.0) -- (1275.0,236.0) ;
\draw [line width=2.0pt, custcol_0_0_0 ]  (1284.0,265.0) -- (1275.0,275.0) -- (1266.0,265.0) -- (1266.0,241.0) -- (1284.0,241.0) -- cycle;
\logisimfontB{\fontsize{12pt}{12pt}\selectfont\node[inner sep=0, outer sep=0, custcol_0_0_0, anchor=base west, rotate=90.0] at  (1282.0,270.0)  {x1};}
\logisimfontC{\fontsize{10pt}{10pt}\fontseries{bx}\selectfont\node[inner sep=0, outer sep=0, custcol_0_0_0, anchor=base west] at  (1255.0,287.0)  {IRWrite};}
\fill [line width=2.0pt, custcol_0_0_0]  (1275.0,235.0) ellipse (2.0 and 2.0 );
\draw [line width=4.0pt, custcol_0_0_0 ]  (475.0,229.0) -- (475.0,228.0) ;
\draw [line width=2.0pt, custcol_0_0_0 ]  (484.0,265.0) -- (475.0,275.0) -- (466.0,265.0) -- (466.0,231.0) -- (484.0,231.0) -- cycle;
\logisimfontB{\fontsize{12pt}{12pt}\selectfont\node[inner sep=0, outer sep=0, custcol_0_0_0, anchor=base west, rotate=90.0] at  (482.0,265.0)  {x2};}
\logisimfontC{\fontsize{10pt}{10pt}\fontseries{bx}\selectfont\node[inner sep=0, outer sep=0, custcol_0_0_0, anchor=base west] at  (451.0,287.0)  {ALUsrcA};}
\fill [line width=2.0pt, custcol_0_0_0]  (475.0,225.0) ellipse (2.0 and 2.0 );
\draw [line width=3.0pt, custcol_0_0_0 ]  (625.0,185.0) -- (625.0,195.0) -- (615.0,195.0) -- (615.0,209.0) ;
\draw [line width=3.0pt, custcol_0_0_0 ]  (575.0,135.0) -- (575.0,195.0) -- (575.0,209.0) ;
\draw [line width=5.0pt, custcol_0_0_0 ]  (615.0,214.0) -- (615.0,209.0) ;
\draw [line width=5.0pt, custcol_0_0_0 ]  (614.0,209.0) -- (576.0,209.0) ;
\logisimfontA{\fontsize{7pt}{7pt}\selectfont\node[inner sep=0, outer sep=0, custcol_0_0_0, anchor=base west, rotate=-90.0] at  (618.0,202.0)  {0};}
\logisimfontA{\fontsize{7pt}{7pt}\selectfont\node[inner sep=0, outer sep=0, custcol_0_0_0, anchor=base west, rotate=-90.0] at  (578.0,202.0)  {1};}
\fill [line width=5.0pt, custcol_0_0_0]  (615.0,215.0) ellipse (2.0 and 2.0 );
\fill [line width=5.0pt, custcol_0_0_0]  (615.0,195.0) ellipse (2.0 and 2.0 );
\fill [line width=5.0pt, custcol_0_0_0]  (575.0,195.0) ellipse (2.0 and 2.0 );
\draw [line width=4.0pt, custcol_0_0_0 ]  (595.0,229.0) -- (595.0,228.0) ;
\draw [line width=2.0pt, custcol_0_0_0 ]  (604.0,265.0) -- (595.0,275.0) -- (586.0,265.0) -- (586.0,231.0) -- (604.0,231.0) -- cycle;
\logisimfontB{\fontsize{12pt}{12pt}\selectfont\node[inner sep=0, outer sep=0, custcol_0_0_0, anchor=base west, rotate=90.0] at  (602.0,265.0)  {x2};}
\logisimfontC{\fontsize{10pt}{10pt}\fontseries{bx}\selectfont\node[inner sep=0, outer sep=0, custcol_0_0_0, anchor=base west] at  (571.0,287.0)  {ALUsrcB};}
\fill [line width=2.0pt, custcol_0_0_0]  (595.0,225.0) ellipse (2.0 and 2.0 );
\draw [line width=3.0pt, custcol_0_0_0 ]  (375.0,135.0) -- (375.0,195.0) -- (375.0,209.0) ;
\draw [line width=3.0pt, custcol_0_0_0 ]  (335.0,135.0) -- (335.0,195.0) -- (335.0,209.0) ;
\draw [line width=5.0pt, custcol_0_0_0 ]  (375.0,214.0) -- (375.0,209.0) ;
\draw [line width=5.0pt, custcol_0_0_0 ]  (374.0,209.0) -- (336.0,209.0) ;
\logisimfontA{\fontsize{7pt}{7pt}\selectfont\node[inner sep=0, outer sep=0, custcol_0_0_0, anchor=base west, rotate=-90.0] at  (378.0,202.0)  {0};}
\logisimfontA{\fontsize{7pt}{7pt}\selectfont\node[inner sep=0, outer sep=0, custcol_0_0_0, anchor=base west, rotate=-90.0] at  (338.0,202.0)  {1};}
\fill [line width=5.0pt, custcol_0_0_0]  (375.0,215.0) ellipse (2.0 and 2.0 );
\fill [line width=5.0pt, custcol_0_0_0]  (375.0,195.0) ellipse (2.0 and 2.0 );
\fill [line width=5.0pt, custcol_0_0_0]  (335.0,195.0) ellipse (2.0 and 2.0 );
\draw [line width=3.0pt, custcol_0_0_0 ]  (495.0,185.0) -- (495.0,195.0) -- (495.0,209.0) ;
\draw [line width=3.0pt, custcol_0_0_0 ]  (455.0,185.0) -- (455.0,195.0) -- (455.0,209.0) ;
\draw [line width=5.0pt, custcol_0_0_0 ]  (495.0,214.0) -- (495.0,209.0) ;
\draw [line width=5.0pt, custcol_0_0_0 ]  (494.0,209.0) -- (456.0,209.0) ;
\logisimfontA{\fontsize{7pt}{7pt}\selectfont\node[inner sep=0, outer sep=0, custcol_0_0_0, anchor=base west, rotate=-90.0] at  (498.0,202.0)  {0};}
\logisimfontA{\fontsize{7pt}{7pt}\selectfont\node[inner sep=0, outer sep=0, custcol_0_0_0, anchor=base west, rotate=-90.0] at  (458.0,202.0)  {1};}
\fill [line width=5.0pt, custcol_0_0_0]  (495.0,215.0) ellipse (2.0 and 2.0 );
\fill [line width=5.0pt, custcol_0_0_0]  (495.0,195.0) ellipse (2.0 and 2.0 );
\fill [line width=5.0pt, custcol_0_0_0]  (455.0,195.0) ellipse (2.0 and 2.0 );
\draw [line width=3.0pt, custcol_0_0_0 ]  (685.0,239.0) -- (685.0,236.0) ;
\draw [line width=2.0pt, custcol_0_0_0 ]  (694.0,265.0) -- (685.0,275.0) -- (676.0,265.0) -- (676.0,241.0) -- (694.0,241.0) -- cycle;
\logisimfontB{\fontsize{12pt}{12pt}\selectfont\node[inner sep=0, outer sep=0, custcol_0_0_0, anchor=base west, rotate=90.0] at  (692.0,270.0)  {x1};}
\logisimfontC{\fontsize{10pt}{10pt}\fontseries{bx}\selectfont\node[inner sep=0, outer sep=0, custcol_0_0_0, anchor=base west] at  (650.0,287.0)  {PCbackWrite};}
\fill [line width=2.0pt, custcol_0_0_0]  (685.0,235.0) ellipse (2.0 and 2.0 );
\draw [line width=4.0pt, custcol_0_0_0 ]  (355.0,229.0) -- (355.0,228.0) ;
\draw [line width=2.0pt, custcol_0_0_0 ]  (364.0,265.0) -- (355.0,275.0) -- (346.0,265.0) -- (346.0,231.0) -- (364.0,231.0) -- cycle;
\logisimfontB{\fontsize{12pt}{12pt}\selectfont\node[inner sep=0, outer sep=0, custcol_0_0_0, anchor=base west, rotate=90.0] at  (362.0,265.0)  {x2};}
\logisimfontC{\fontsize{10pt}{10pt}\fontseries{bx}\selectfont\node[inner sep=0, outer sep=0, custcol_0_0_0, anchor=base west] at  (336.0,287.0)  {ALUOp};}
\fill [line width=2.0pt, custcol_0_0_0]  (355.0,225.0) ellipse (2.0 and 2.0 );
\draw [line width=3.0pt, custcol_0_0_0 ]  (275.0,239.0) -- (275.0,236.0) ;
\draw [line width=2.0pt, custcol_0_0_0 ]  (284.0,265.0) -- (275.0,275.0) -- (266.0,265.0) -- (266.0,241.0) -- (284.0,241.0) -- cycle;
\logisimfontB{\fontsize{12pt}{12pt}\selectfont\node[inner sep=0, outer sep=0, custcol_0_0_0, anchor=base west, rotate=90.0] at  (282.0,270.0)  {x1};}
\logisimfontC{\fontsize{10pt}{10pt}\fontseries{bx}\selectfont\node[inner sep=0, outer sep=0, custcol_0_0_0, anchor=base west] at  (258.0,287.0)  {PCsrc};}
\fill [line width=2.0pt, custcol_0_0_0]  (275.0,235.0) ellipse (2.0 and 2.0 );
\end{tikzpicture}

}
    \vspace{-20pt}
    \caption{Esquemático da lógica de saída do controle}
\end{figure}

\begin{table}[H]
    \centering
    \resizebox{\textwidth}{!}{
    \begin{tabular}{|c|c|c|c|c|c|c|c|c|}
        \hline
        \rowcolor{black}
        \textcolor{white}{\textbf{Estado}} & \multicolumn{8}{|c|}{\textcolor{white}{\textbf{Próximo Estado}}} \\ \cline{2-9}
        \rowcolor{black}
        \textcolor{white}{\textbf{Atual}} & \textcolor{white}{\textbf{\scriptsize 0000011}} &
        \textcolor{white}{\textbf{\scriptsize 0100011}} &
        \textcolor{white}{\textbf{\scriptsize 0010011}} &
        \textcolor{white}{\textbf{\scriptsize 1100111}} &
        \textcolor{white}{\textbf{\scriptsize 0110011}} &
        \textcolor{white}{\textbf{\scriptsize 1100011}} &
        \textcolor{white}{\textbf{\scriptsize 1101111}} &
        \textcolor{white}{\textbf{\scriptsize 0110111}} \\ \hline
         
        0000 & 0001 & 0001 & 0001 & 0001 & 0001 & 0001 & 0001 & 0001 \\ \hline

        \rowcolor{gray!20}
        0001 & 0011 & 0011 & 0011 & 0011 & 0010 & 1000 & 0101 & 1101 \\ \hline

        0011 & 0110 & 0111 & 1001 & 0101 & X & X & X & X \\ \hline

        \rowcolor{gray!20}
        0010 & X & X & X & X & 1001 & X & X & X \\ \hline
        
        1000 & X & X & X & X & X & 0000 & X & X \\ \hline

        \rowcolor{gray!20}
        0101 & X & X & X & 0000 & X & X & 0000 & X \\ \hline

        0110 & 1011 & X & X & X & X & X & X & X \\ \hline

        \rowcolor{gray!20}
        0111 & X & 0000 & X & X & X & X & X & X \\ \hline

        1001 & X & X & 0000 & X & 0000 & X & X & X \\ \hline

        \rowcolor{gray!20}
        1011 & X & 0000 & X & X & X & X & X & X \\ \hline

        1101 & X & X & X & X & X & X & X & 0000 \\ \hline
    \end{tabular}
    }
    \caption{Tabela de transição de estados do controle}
    \label{tab:tran_estados}
\end{table}

A abordagem para implementar a lógica de próximo estado é um pouco diferente; não minimizaremos equações nem nada do tipo. Na verdade, tendo definido que um registrador de deslocamento será utilizado como elemento de memória, montamos outra tabela a partir da \autoref{tab:tran_estados}; uma em que, ao invés do próximo estado, representamos, em cada célula, a operação que o registrador precisa fazer para chegar àquele estado.

O Registrador de Deslocamento é capaz de fazer quatro operações: \textit{shift left(L)}, \textit{shift left(R)}, \textit{load(D)} e \textit{reset}, em que a variável entre parenteses representa o argumento necessário para aquela operação (L: o bit que entra a esquerda no \textit{shift left}, R: o bit que entra a direita \textit{shift left}, D: o vetor de quatro bits que substitui o estado atual no \textit{load}). O \textit{reset} não foi utilizado nesse projeto, ao invés dele, é utilizado o \textit{load(0000)}.

\begin{figure}[H]
    \centering
    % Important: If latex complains about unicode characters,
% please use "\usepackage[utf8x]{inputenc}" in your preamble
% You can change the size of the picture by putting it into the construct:
% 1) \resizebox{10cm}{!}{"below picture"} to scale horizontally to 10 cm
% 2) \resizebox{!}{15cm}{"below picture"} to scale vertically to 15 cm
% 3) \resizebox{10cm}{15cm}{"below picture"} a combination of above two
% It is not recomended to use the scale option of the tikzpicture environment.
\resizebox{0.5\textwidth}{!}{
\begin{tikzpicture}[x=1pt,y=-1pt,line cap=rect]
\def\logisimfontA#1{\fontfamily{cmr}{#1}} % Replaced by logisim, original font was "SansSerif"
\def\logisimfontB#1{\fontfamily{Ubuntu}{#1}}
\def\logisimfontC#1{\fontfamily{cmtt}{#1}} % Replaced by logisim, original font was "Monospaced"
\def\logisimfontD#1{\fontfamily{Arial Rounded MT Bold}{#1}}
\definecolor{custcol_0_0_0}{RGB}{0, 0, 0}
\definecolor{custcol_75_75_75}{RGB}{117, 117, 117}
\definecolor{custcol_ff_ff_ff}{RGB}{255, 255, 255}
\draw [line width=3.0pt, custcol_0_0_0 ]  (345.0,538.0) -- (345.0,548.0) ;
\draw [line width=4.0pt, custcol_0_0_0 ]  (255.0,128.0) -- (255.0,148.0) ;
\draw [line width=3.0pt, custcol_0_0_0 ]  (265.0,678.0) -- (295.0,678.0) ;
\draw [line width=3.0pt, custcol_0_0_0 ]  (185.0,58.0) -- (185.0,98.0) -- (235.0,98.0) ;
\draw [line width=3.0pt, custcol_0_0_0 ]  (55.0,58.0) -- (55.0,68.0) ;
\draw [line width=3.0pt, custcol_0_0_0 ]  (95.0,58.0) -- (95.0,68.0) ;
\draw [line width=4.0pt, custcol_0_0_0 ]  (15.0,58.0) -- (15.0,68.0) ;
\draw [line width=3.0pt, custcol_0_0_0 ]  (325.0,58.0) -- (325.0,48.0) -- (365.0,48.0) -- (365.0,58.0) ;
\draw [line width=3.0pt, custcol_0_0_0 ]  (345.0,348.0) -- (345.0,358.0) ;
\draw [line width=4.0pt, custcol_0_0_0 ]  (255.0,698.0) -- (255.0,718.0) ;
\draw [line width=4.0pt, custcol_0_0_0 ]  (135.0,58.0) -- (135.0,88.0) ;
\draw [line width=4.0pt, custcol_0_0_0 ]  (585.0,98.0) -- (605.0,98.0) ;
\draw [line width=3.0pt, custcol_0_0_0 ]  (265.0,488.0) -- (295.0,488.0) ;
\draw [line width=3.0pt, custcol_0_0_0 ]  (345.0,158.0) -- (345.0,168.0) ;
\draw [line width=3.0pt, custcol_0_0_0 ]  (365.0,608.0) -- (365.0,618.0) -- (325.0,618.0) -- (325.0,628.0) ;
\draw [line width=3.0pt, custcol_0_0_0 ]  (325.0,418.0) -- (325.0,428.0) -- (205.0,428.0) -- (205.0,308.0) -- (235.0,308.0) ;
\draw [line width=3.0pt, custcol_0_0_0 ]  (205.0,238.0) -- (325.0,238.0) -- (325.0,248.0) ;
\draw [line width=3.0pt, custcol_0_0_0 ]  (365.0,418.0) -- (365.0,428.0) -- (425.0,428.0) -- (425.0,488.0) -- (395.0,488.0) ;
\draw [line width=4.0pt, custcol_0_0_0 ]  (255.0,508.0) -- (255.0,528.0) ;
\draw [line width=3.0pt, custcol_0_0_0 ]  (395.0,108.0) -- (425.0,108.0) -- (475.0,108.0) -- (475.0,198.0) -- (195.0,198.0) -- (195.0,288.0) -- (235.0,288.0) ;
\draw [line width=3.0pt, custcol_0_0_0 ]  (235.0,278.0) -- (205.0,278.0) -- (205.0,238.0) -- (205.0,118.0) -- (235.0,118.0) ;
\draw [line width=3.0pt, custcol_0_0_0 ]  (265.0,298.0) -- (295.0,298.0) ;
\draw [line width=3.0pt, custcol_0_0_0 ]  (325.0,608.0) -- (325.0,618.0) -- (205.0,618.0) -- (205.0,658.0) -- (235.0,658.0) ;
\draw [line width=3.0pt, custcol_0_0_0 ]  (365.0,248.0) -- (365.0,238.0) -- (425.0,238.0) -- (425.0,298.0) ;
\draw [line width=3.0pt, custcol_0_0_0 ]  (325.0,438.0) -- (325.0,428.0) -- (365.0,428.0) -- (365.0,438.0) ;
\draw [line width=3.0pt, custcol_0_0_0 ]  (345.0,728.0) -- (345.0,738.0) ;
\draw [line width=4.0pt, custcol_0_0_0 ]  (255.0,318.0) -- (255.0,338.0) ;
\draw [line width=3.0pt, custcol_0_0_0 ]  (475.0,488.0) -- (475.0,578.0) -- (195.0,578.0) -- (195.0,668.0) -- (235.0,668.0) ;
\draw [line width=3.0pt, custcol_0_0_0 ]  (325.0,228.0) -- (325.0,238.0) -- (365.0,238.0) -- (365.0,228.0) ;
\draw [line width=3.0pt, custcol_0_0_0 ]  (395.0,678.0) -- (425.0,678.0) ;
\draw [line width=3.0pt, custcol_0_0_0 ]  (325.0,28.0) -- (325.0,48.0) -- (205.0,48.0) -- (205.0,88.0) -- (235.0,88.0) ;
\draw [line width=3.0pt, custcol_0_0_0 ]  (265.0,108.0) -- (295.0,108.0) ;
\draw [line width=3.0pt, custcol_0_0_0 ]  (205.0,428.0) -- (205.0,468.0) -- (235.0,468.0) ;
\draw [line width=3.0pt, custcol_0_0_0 ]  (235.0,478.0) -- (195.0,478.0) -- (195.0,388.0) -- (475.0,388.0) -- (475.0,298.0) -- (425.0,298.0) -- (395.0,298.0) ;
\draw [line width=3.0pt, custcol_0_0_0 ]  (165.0,58.0) -- (165.0,688.0) -- (235.0,688.0) ;
\draw [line width=3.0pt, custcol_0_0_0 ]  (365.0,28.0) -- (365.0,48.0) -- (425.0,48.0) -- (425.0,108.0) ;
\draw [line width=3.0pt, custcol_0_0_0 ]  (235.0,498.0) -- (205.0,498.0) -- (205.0,618.0) ;
\fill [line width=3.0pt, custcol_0_0_0]  (475.0,488.0) ellipse (5.0 and 5.0 );
\fill [line width=3.0pt, custcol_0_0_0]  (205.0,428.0) ellipse (5.0 and 5.0 );
\fill [line width=3.0pt, custcol_0_0_0]  (325.0,48.0) ellipse (5.0 and 5.0 );
\fill [line width=3.0pt, custcol_0_0_0]  (365.0,618.0) ellipse (5.0 and 5.0 );
\fill [line width=3.0pt, custcol_0_0_0]  (425.0,298.0) ellipse (5.0 and 5.0 );
\fill [line width=3.0pt, custcol_0_0_0]  (475.0,298.0) ellipse (5.0 and 5.0 );
\fill [line width=3.0pt, custcol_0_0_0]  (205.0,238.0) ellipse (5.0 and 5.0 );
\fill [line width=3.0pt, custcol_0_0_0]  (325.0,618.0) ellipse (5.0 and 5.0 );
\fill [line width=3.0pt, custcol_0_0_0]  (365.0,428.0) ellipse (5.0 and 5.0 );
\fill [line width=3.0pt, custcol_0_0_0]  (425.0,108.0) ellipse (5.0 and 5.0 );
\fill [line width=3.0pt, custcol_0_0_0]  (475.0,108.0) ellipse (5.0 and 5.0 );
\fill [line width=3.0pt, custcol_0_0_0]  (325.0,428.0) ellipse (5.0 and 5.0 );
\fill [line width=3.0pt, custcol_0_0_0]  (425.0,678.0) ellipse (5.0 and 5.0 );
\fill [line width=3.0pt, custcol_0_0_0]  (365.0,238.0) ellipse (5.0 and 5.0 );
\fill [line width=3.0pt, custcol_0_0_0]  (205.0,618.0) ellipse (5.0 and 5.0 );
\fill [line width=3.0pt, custcol_0_0_0]  (325.0,238.0) ellipse (5.0 and 5.0 );
\fill [line width=3.0pt, custcol_0_0_0]  (425.0,488.0) ellipse (5.0 and 5.0 );
\fill [line width=3.0pt, custcol_0_0_0]  (365.0,48.0) ellipse (5.0 and 5.0 );
\draw [line width=2.0pt, custcol_0_0_0 ]  (7.0,40.0) -- (24.0,40.0) ;
\draw [line width=2.0pt, custcol_0_0_0 ]  (25.0,40.0) -- (25.0,57.0) ;
\draw [line width=2.0pt, custcol_0_0_0 ]  (25.0,58.0) -- (8.0,58.0) ;
\draw [line width=2.0pt, custcol_0_0_0 ]  (7.0,58.0) -- (7.0,41.0) ;
\logisimfontA{\fontsize{12pt}{12pt}\selectfont\node[inner sep=0, outer sep=0, custcol_0_0_0, anchor=base west] at  (9.0,55.0)  {x2};}
\logisimfontA{\fontsize{16pt}{16pt}\fontseries{bx}\selectfont\node[inner sep=0, outer sep=0, custcol_0_0_0, anchor=base west] at  (10.0,32.0)  {S};}
\fill [line width=2.0pt, custcol_0_0_0]  (15.0,58.0) ellipse (2.0 and 2.0 );
\logisimfontB{\fontsize{16pt}{16pt}\fontseries{bx}\selectfont\node[inner sep=0, outer sep=0, custcol_0_0_0, anchor=base west] at  (11.0,87.0)  {S};}
\draw [line width=2.0pt, custcol_0_0_0 ]  (7.0,72.0) -- (15.0,68.0) -- (23.0,72.0) -- (23.0,94.0) -- (7.0,94.0) -- cycle;
\fill [line width=2.0pt, custcol_0_0_0]  (15.0,68.0) ellipse (2.0 and 2.0 );
\draw [line width=2.0pt, custcol_0_0_0 ]  (47.0,40.0) -- (64.0,40.0) ;
\draw [line width=2.0pt, custcol_0_0_0 ]  (65.0,40.0) -- (65.0,57.0) ;
\draw [line width=2.0pt, custcol_0_0_0 ]  (65.0,58.0) -- (48.0,58.0) ;
\draw [line width=2.0pt, custcol_0_0_0 ]  (47.0,58.0) -- (47.0,41.0) ;
\logisimfontA{\fontsize{12pt}{12pt}\selectfont\node[inner sep=0, outer sep=0, custcol_0_0_0, anchor=base west] at  (49.0,55.0)  {x1};}
\logisimfontA{\fontsize{16pt}{16pt}\fontseries{bx}\selectfont\node[inner sep=0, outer sep=0, custcol_0_0_0, anchor=base west] at  (39.0,32.0)  {CLK};}
\fill [line width=2.0pt, custcol_0_0_0]  (55.0,58.0) ellipse (2.0 and 2.0 );
\logisimfontB{\fontsize{12pt}{12pt}\fontseries{bx}\selectfont\node[inner sep=0, outer sep=0, custcol_0_0_0, anchor=base west] at  (44.0,84.0)  {CLK};}
\draw [line width=2.0pt, custcol_0_0_0 ]  (41.0,72.0) -- (45.0,72.0) -- (55.0,68.0) -- (65.0,72.0) -- (69.0,72.0) -- (69.0,91.0) -- (41.0,91.0) -- cycle;
\fill [line width=2.0pt, custcol_0_0_0]  (55.0,68.0) ellipse (2.0 and 2.0 );
\draw [line width=2.0pt, custcol_0_0_0 ]  (87.0,40.0) -- (104.0,40.0) ;
\draw [line width=2.0pt, custcol_0_0_0 ]  (105.0,40.0) -- (105.0,57.0) ;
\draw [line width=2.0pt, custcol_0_0_0 ]  (105.0,58.0) -- (88.0,58.0) ;
\draw [line width=2.0pt, custcol_0_0_0 ]  (87.0,58.0) -- (87.0,41.0) ;
\logisimfontA{\fontsize{12pt}{12pt}\selectfont\node[inner sep=0, outer sep=0, custcol_0_0_0, anchor=base west] at  (89.0,55.0)  {x1};}
\logisimfontB{\fontsize{16pt}{16pt}\fontseries{bx}\selectfont\node[inner sep=0, outer sep=0, custcol_0_0_0, anchor=base west] at  (81.0,33.0)  {RST};}
\fill [line width=2.0pt, custcol_0_0_0]  (95.0,58.0) ellipse (2.0 and 2.0 );
\logisimfontB{\fontsize{12pt}{12pt}\fontseries{bx}\selectfont\node[inner sep=0, outer sep=0, custcol_0_0_0, anchor=base west] at  (85.0,84.0)  {RST};}
\draw [line width=2.0pt, custcol_0_0_0 ]  (81.0,72.0) -- (85.0,72.0) -- (95.0,68.0) -- (105.0,72.0) -- (109.0,72.0) -- (109.0,91.0) -- (81.0,91.0) -- cycle;
\fill [line width=2.0pt, custcol_0_0_0]  (95.0,68.0) ellipse (2.0 and 2.0 );
\draw [line width=3.0pt, custcol_0_0_0 ]  (235.0,128.0) -- (145.0,128.0) -- (145.0,108.0) -- (145.0,94.0) ;
\draw [line width=3.0pt, custcol_0_0_0 ]  (235.0,318.0) -- (135.0,318.0) -- (135.0,108.0) -- (135.0,94.0) ;
\draw [line width=3.0pt, custcol_0_0_0 ]  (235.0,508.0) -- (125.0,508.0) -- (125.0,108.0) -- (125.0,94.0) ;
\draw [line width=3.0pt, custcol_0_0_0 ]  (235.0,698.0) -- (115.0,698.0) -- (115.0,108.0) -- (115.0,94.0) ;
\draw [line width=5.0pt, custcol_0_0_0 ]  (135.0,89.0) -- (135.0,94.0) ;
\draw [line width=5.0pt, custcol_0_0_0 ]  (144.0,94.0) -- (116.0,94.0) ;
\logisimfontA{\fontsize{7pt}{7pt}\selectfont\node[inner sep=0, outer sep=0, custcol_0_0_0, anchor=base west, rotate=-90.0] at  (148.0,97.0)  {3};}
\logisimfontA{\fontsize{7pt}{7pt}\selectfont\node[inner sep=0, outer sep=0, custcol_0_0_0, anchor=base west, rotate=-90.0] at  (138.0,97.0)  {2};}
\logisimfontA{\fontsize{7pt}{7pt}\selectfont\node[inner sep=0, outer sep=0, custcol_0_0_0, anchor=base west, rotate=-90.0] at  (128.0,97.0)  {1};}
\logisimfontA{\fontsize{7pt}{7pt}\selectfont\node[inner sep=0, outer sep=0, custcol_0_0_0, anchor=base west, rotate=-90.0] at  (118.0,97.0)  {0};}
\fill [line width=5.0pt, custcol_0_0_0]  (135.0,88.0) ellipse (2.0 and 2.0 );
\fill [line width=5.0pt, custcol_0_0_0]  (145.0,108.0) ellipse (2.0 and 2.0 );
\fill [line width=5.0pt, custcol_0_0_0]  (135.0,108.0) ellipse (2.0 and 2.0 );
\fill [line width=5.0pt, custcol_0_0_0]  (125.0,108.0) ellipse (2.0 and 2.0 );
\fill [line width=5.0pt, custcol_0_0_0]  (115.0,108.0) ellipse (2.0 and 2.0 );
\draw [line width=2.0pt, custcol_0_0_0 ]  (117.0,40.0) -- (154.0,40.0) ;
\draw [line width=2.0pt, custcol_0_0_0 ]  (155.0,40.0) -- (155.0,57.0) ;
\draw [line width=2.0pt, custcol_0_0_0 ]  (155.0,58.0) -- (118.0,58.0) ;
\draw [line width=2.0pt, custcol_0_0_0 ]  (117.0,58.0) -- (117.0,41.0) ;
\logisimfontA{\fontsize{12pt}{12pt}\selectfont\node[inner sep=0, outer sep=0, custcol_0_0_0, anchor=base west] at  (129.0,55.0)  {x4};}
\logisimfontB{\fontsize{16pt}{16pt}\fontseries{bx}\selectfont\node[inner sep=0, outer sep=0, custcol_0_0_0, anchor=base west] at  (129.0,33.0)  {D};}
\fill [line width=2.0pt, custcol_0_0_0]  (135.0,58.0) ellipse (2.0 and 2.0 );
\draw [line width=2.0pt, custcol_0_0_0 ]  (157.0,40.0) -- (174.0,40.0) ;
\draw [line width=2.0pt, custcol_0_0_0 ]  (175.0,40.0) -- (175.0,57.0) ;
\draw [line width=2.0pt, custcol_0_0_0 ]  (175.0,58.0) -- (158.0,58.0) ;
\draw [line width=2.0pt, custcol_0_0_0 ]  (157.0,58.0) -- (157.0,41.0) ;
\logisimfontA{\fontsize{12pt}{12pt}\selectfont\node[inner sep=0, outer sep=0, custcol_0_0_0, anchor=base west] at  (159.0,55.0)  {x1};}
\logisimfontB{\fontsize{16pt}{16pt}\fontseries{bx}\selectfont\node[inner sep=0, outer sep=0, custcol_0_0_0, anchor=base west] at  (161.0,33.0)  {L};}
\fill [line width=2.0pt, custcol_0_0_0]  (165.0,58.0) ellipse (2.0 and 2.0 );
\draw [line width=2.0pt, custcol_0_0_0 ]  (177.0,40.0) -- (194.0,40.0) ;
\draw [line width=2.0pt, custcol_0_0_0 ]  (195.0,40.0) -- (195.0,57.0) ;
\draw [line width=2.0pt, custcol_0_0_0 ]  (195.0,58.0) -- (178.0,58.0) ;
\draw [line width=2.0pt, custcol_0_0_0 ]  (177.0,58.0) -- (177.0,41.0) ;
\logisimfontA{\fontsize{12pt}{12pt}\selectfont\node[inner sep=0, outer sep=0, custcol_0_0_0, anchor=base west] at  (179.0,55.0)  {x1};}
\logisimfontA{\fontsize{16pt}{16pt}\fontseries{bx}\selectfont\node[inner sep=0, outer sep=0, custcol_0_0_0, anchor=base west] at  (180.0,32.0)  {R};}
\fill [line width=2.0pt, custcol_0_0_0]  (185.0,58.0) ellipse (2.0 and 2.0 );
\logisimfontB{\fontsize{16pt}{16pt}\fontseries{bx}\selectfont\node[inner sep=0, outer sep=0, custcol_0_0_0, anchor=base west] at  (251.0,167.0)  {S};}
\draw [line width=2.0pt, custcol_0_0_0 ]  (247.0,152.0) -- (255.0,148.0) -- (263.0,152.0) -- (263.0,174.0) -- (247.0,174.0) -- cycle;
\fill [line width=2.0pt, custcol_0_0_0]  (255.0,148.0) ellipse (2.0 and 2.0 );
\logisimfontB{\fontsize{16pt}{16pt}\fontseries{bx}\selectfont\node[inner sep=0, outer sep=0, custcol_0_0_0, anchor=base west] at  (251.0,357.0)  {S};}
\draw [line width=2.0pt, custcol_0_0_0 ]  (247.0,342.0) -- (255.0,338.0) -- (263.0,342.0) -- (263.0,364.0) -- (247.0,364.0) -- cycle;
\fill [line width=2.0pt, custcol_0_0_0]  (255.0,338.0) ellipse (2.0 and 2.0 );
\logisimfontB{\fontsize{16pt}{16pt}\fontseries{bx}\selectfont\node[inner sep=0, outer sep=0, custcol_0_0_0, anchor=base west] at  (251.0,547.0)  {S};}
\draw [line width=2.0pt, custcol_0_0_0 ]  (247.0,532.0) -- (255.0,528.0) -- (263.0,532.0) -- (263.0,554.0) -- (247.0,554.0) -- cycle;
\fill [line width=2.0pt, custcol_0_0_0]  (255.0,528.0) ellipse (2.0 and 2.0 );
\logisimfontB{\fontsize{16pt}{16pt}\fontseries{bx}\selectfont\node[inner sep=0, outer sep=0, custcol_0_0_0, anchor=base west] at  (251.0,737.0)  {S};}
\draw [line width=2.0pt, custcol_0_0_0 ]  (247.0,722.0) -- (255.0,718.0) -- (263.0,722.0) -- (263.0,744.0) -- (247.0,744.0) -- cycle;
\fill [line width=2.0pt, custcol_0_0_0]  (255.0,718.0) ellipse (2.0 and 2.0 );
\draw [line width=2.0pt, custcol_0_0_0 ]  (235.0,78.0) -- (235.0,138.0) -- (264.0,123.0) -- (264.0,93.0) -- cycle;
\fill [line width=1.0pt, custcol_0_0_0]  (235.0,88.0) ellipse (2.0 and 2.0 );
\fill [line width=1.0pt, custcol_0_0_0]  (235.0,98.0) ellipse (2.0 and 2.0 );
\fill [line width=1.0pt, custcol_0_0_0]  (235.0,118.0) ellipse (2.0 and 2.0 );
\fill [line width=1.0pt, custcol_0_0_0]  (235.0,128.0) ellipse (2.0 and 2.0 );
\fill [line width=1.0pt, custcol_0_0_0]  (255.0,128.0) ellipse (2.0 and 2.0 );
\fill [line width=1.0pt, custcol_0_0_0]  (265.0,108.0) ellipse (2.0 and 2.0 );
\draw [line width=2.0pt, custcol_0_0_0 ]  (235.0,268.0) -- (235.0,328.0) -- (264.0,313.0) -- (264.0,283.0) -- cycle;
\fill [line width=1.0pt, custcol_0_0_0]  (235.0,278.0) ellipse (2.0 and 2.0 );
\fill [line width=1.0pt, custcol_0_0_0]  (235.0,288.0) ellipse (2.0 and 2.0 );
\fill [line width=1.0pt, custcol_0_0_0]  (235.0,308.0) ellipse (2.0 and 2.0 );
\fill [line width=1.0pt, custcol_0_0_0]  (235.0,318.0) ellipse (2.0 and 2.0 );
\fill [line width=1.0pt, custcol_0_0_0]  (255.0,318.0) ellipse (2.0 and 2.0 );
\fill [line width=1.0pt, custcol_0_0_0]  (265.0,298.0) ellipse (2.0 and 2.0 );
\draw [line width=2.0pt, custcol_0_0_0 ]  (235.0,458.0) -- (235.0,518.0) -- (264.0,503.0) -- (264.0,473.0) -- cycle;
\fill [line width=1.0pt, custcol_0_0_0]  (235.0,468.0) ellipse (2.0 and 2.0 );
\fill [line width=1.0pt, custcol_0_0_0]  (235.0,478.0) ellipse (2.0 and 2.0 );
\fill [line width=1.0pt, custcol_0_0_0]  (235.0,498.0) ellipse (2.0 and 2.0 );
\fill [line width=1.0pt, custcol_0_0_0]  (235.0,508.0) ellipse (2.0 and 2.0 );
\fill [line width=1.0pt, custcol_0_0_0]  (255.0,508.0) ellipse (2.0 and 2.0 );
\fill [line width=1.0pt, custcol_0_0_0]  (265.0,488.0) ellipse (2.0 and 2.0 );
\draw [line width=2.0pt, custcol_0_0_0 ]  (235.0,648.0) -- (235.0,708.0) -- (264.0,693.0) -- (264.0,663.0) -- cycle;
\fill [line width=1.0pt, custcol_0_0_0]  (235.0,658.0) ellipse (2.0 and 2.0 );
\fill [line width=1.0pt, custcol_0_0_0]  (235.0,668.0) ellipse (2.0 and 2.0 );
\fill [line width=1.0pt, custcol_0_0_0]  (235.0,688.0) ellipse (2.0 and 2.0 );
\fill [line width=1.0pt, custcol_0_0_0]  (235.0,698.0) ellipse (2.0 and 2.0 );
\fill [line width=1.0pt, custcol_0_0_0]  (255.0,698.0) ellipse (2.0 and 2.0 );
\fill [line width=1.0pt, custcol_0_0_0]  (265.0,678.0) ellipse (2.0 and 2.0 );
\logisimfontC{\fontsize{12pt}{12pt}\selectfont\node[inner sep=0, outer sep=0, custcol_0_0_0, anchor=base west] at  (322.0,225.0)  {0};}
\fill [line width=1.0pt, custcol_0_0_0]  (325.0,228.0) ellipse (2.0 and 2.0 );
\logisimfontC{\fontsize{12pt}{12pt}\selectfont\node[inner sep=0, outer sep=0, custcol_0_0_0, anchor=base west] at  (322.0,25.0)  {0};}
\fill [line width=1.0pt, custcol_0_0_0]  (325.0,28.0) ellipse (2.0 and 2.0 );
\logisimfontC{\fontsize{12pt}{12pt}\selectfont\node[inner sep=0, outer sep=0, custcol_0_0_0, anchor=base west] at  (322.0,415.0)  {0};}
\fill [line width=1.0pt, custcol_0_0_0]  (325.0,418.0) ellipse (2.0 and 2.0 );
\logisimfontB{\fontsize{12pt}{12pt}\fontseries{bx}\selectfont\node[inner sep=0, outer sep=0, custcol_0_0_0, anchor=base west] at  (334.0,184.0)  {CLK};}
\draw [line width=2.0pt, custcol_0_0_0 ]  (331.0,172.0) -- (335.0,172.0) -- (345.0,168.0) -- (355.0,172.0) -- (359.0,172.0) -- (359.0,191.0) -- (331.0,191.0) -- cycle;
\fill [line width=2.0pt, custcol_0_0_0]  (345.0,168.0) ellipse (2.0 and 2.0 );
\begin{pgfpicture}
   \begin{pgfmagnify}{1pt}{-1pt}
      \pgfsetrectcap
      \pgfsetcornersarced{\pgfpoint{3.0}{3.0}}
      \pgfsetlinewidth{2.0}
      \color{custcol_0_0_0}
      \pgfsetfillopacity{1.0}
      \pgfpathrectanglecorners{\pgfpoint{607.0}{90.0}}{\pgfpoint{645.0}{108.0}}
      \pgfusepath{stroke}
   \end{pgfmagnify}
\end{pgfpicture}
\logisimfontA{\fontsize{12pt}{12pt}\selectfont\node[inner sep=0, outer sep=0, custcol_0_0_0, anchor=base west] at  (619.0,105.0)  {x4};}
\logisimfontB{\fontsize{16pt}{16pt}\fontseries{bx}\selectfont\node[inner sep=0, outer sep=0, custcol_0_0_0, anchor=base west] at  (647.0,105.0)  {Q};}
\fill [line width=2.0pt, custcol_0_0_0]  (605.0,98.0) ellipse (2.0 and 2.0 );
\logisimfontB{\fontsize{12pt}{12pt}\fontseries{bx}\selectfont\node[inner sep=0, outer sep=0, custcol_0_0_0, anchor=base west] at  (334.0,754.0)  {CLK};}
\draw [line width=2.0pt, custcol_0_0_0 ]  (331.0,742.0) -- (335.0,742.0) -- (345.0,738.0) -- (355.0,742.0) -- (359.0,742.0) -- (359.0,761.0) -- (331.0,761.0) -- cycle;
\fill [line width=2.0pt, custcol_0_0_0]  (345.0,738.0) ellipse (2.0 and 2.0 );
\logisimfontB{\fontsize{12pt}{12pt}\fontseries{bx}\selectfont\node[inner sep=0, outer sep=0, custcol_0_0_0, anchor=base west] at  (355.0,409.0)  {RST};}
\draw [line width=2.0pt, custcol_0_0_0 ]  (351.0,395.0) -- (379.0,395.0) -- (379.0,415.0) -- (375.0,415.0) -- (365.0,419.0) -- (355.0,415.0) -- (351.0,415.0) -- cycle;
\fill [line width=2.0pt, custcol_0_0_0]  (365.0,418.0) ellipse (2.0 and 2.0 );
\draw [line width=2.0pt, custcol_0_0_0 ]  (295.0,438.0) -- (394.0,438.0) ;
\draw [line width=2.0pt, custcol_0_0_0 ]  (395.0,438.0) -- (395.0,537.0) ;
\draw [line width=2.0pt, custcol_0_0_0 ]  (395.0,538.0) -- (296.0,538.0) ;
\draw [line width=2.0pt, custcol_0_0_0 ]  (295.0,538.0) -- (295.0,439.0) ;
\draw [line width=1.0pt, custcol_0_0_0 ]  (345.0,528.0) -- (343.0,538.0) -- (347.0,538.0) -- cycle;
\logisimfontD{\fontsize{14pt}{14pt}\fontseries{bx}\selectfont\node[inner sep=0, outer sep=0, custcol_0_0_0, anchor=base west] at  (375.0,493.0)  {Q};}
\logisimfontD{\fontsize{14pt}{14pt}\fontseries{bx}\selectfont\node[inner sep=0, outer sep=0, custcol_75_75_75, anchor=base west] at  (315.0,456.0)  {PR};}
\logisimfontD{\fontsize{14pt}{14pt}\fontseries{bx}\selectfont\node[inner sep=0, outer sep=0, custcol_75_75_75, anchor=base west] at  (351.0,456.0)  {CLR};}
\logisimfontD{\fontsize{14pt}{14pt}\fontseries{bx}\selectfont\node[inner sep=0, outer sep=0, custcol_0_0_0, anchor=base west] at  (303.0,493.0)  {D};}
\fill [line width=1.0pt, custcol_0_0_0]  (295.0,488.0) ellipse (2.0 and 2.0 );
\fill [line width=1.0pt, custcol_0_0_0]  (325.0,438.0) ellipse (2.0 and 2.0 );
\fill [line width=1.0pt, custcol_0_0_0]  (345.0,538.0) ellipse (2.0 and 2.0 );
\fill [line width=1.0pt, custcol_0_0_0]  (365.0,438.0) ellipse (2.0 and 2.0 );
\fill [line width=1.0pt, custcol_0_0_0]  (395.0,488.0) ellipse (2.0 and 2.0 );
\logisimfontC{\fontsize{12pt}{12pt}\selectfont\node[inner sep=0, outer sep=0, custcol_0_0_0, anchor=base west] at  (322.0,605.0)  {0};}
\fill [line width=1.0pt, custcol_0_0_0]  (325.0,608.0) ellipse (2.0 and 2.0 );
\logisimfontB{\fontsize{12pt}{12pt}\fontseries{bx}\selectfont\node[inner sep=0, outer sep=0, custcol_0_0_0, anchor=base west] at  (334.0,374.0)  {CLK};}
\draw [line width=2.0pt, custcol_0_0_0 ]  (331.0,362.0) -- (335.0,362.0) -- (345.0,358.0) -- (355.0,362.0) -- (359.0,362.0) -- (359.0,381.0) -- (331.0,381.0) -- cycle;
\fill [line width=2.0pt, custcol_0_0_0]  (345.0,358.0) ellipse (2.0 and 2.0 );
\draw [line width=2.0pt, custcol_0_0_0 ]  (295.0,628.0) -- (394.0,628.0) ;
\draw [line width=2.0pt, custcol_0_0_0 ]  (395.0,628.0) -- (395.0,727.0) ;
\draw [line width=2.0pt, custcol_0_0_0 ]  (395.0,728.0) -- (296.0,728.0) ;
\draw [line width=2.0pt, custcol_0_0_0 ]  (295.0,728.0) -- (295.0,629.0) ;
\draw [line width=1.0pt, custcol_0_0_0 ]  (345.0,718.0) -- (343.0,728.0) -- (347.0,728.0) -- cycle;
\logisimfontD{\fontsize{14pt}{14pt}\fontseries{bx}\selectfont\node[inner sep=0, outer sep=0, custcol_0_0_0, anchor=base west] at  (375.0,683.0)  {Q};}
\logisimfontD{\fontsize{14pt}{14pt}\fontseries{bx}\selectfont\node[inner sep=0, outer sep=0, custcol_75_75_75, anchor=base west] at  (315.0,646.0)  {PR};}
\logisimfontD{\fontsize{14pt}{14pt}\fontseries{bx}\selectfont\node[inner sep=0, outer sep=0, custcol_75_75_75, anchor=base west] at  (351.0,646.0)  {CLR};}
\logisimfontD{\fontsize{14pt}{14pt}\fontseries{bx}\selectfont\node[inner sep=0, outer sep=0, custcol_0_0_0, anchor=base west] at  (303.0,683.0)  {D};}
\fill [line width=1.0pt, custcol_0_0_0]  (295.0,678.0) ellipse (2.0 and 2.0 );
\fill [line width=1.0pt, custcol_0_0_0]  (325.0,628.0) ellipse (2.0 and 2.0 );
\fill [line width=1.0pt, custcol_0_0_0]  (345.0,728.0) ellipse (2.0 and 2.0 );
\fill [line width=1.0pt, custcol_0_0_0]  (365.0,628.0) ellipse (2.0 and 2.0 );
\fill [line width=1.0pt, custcol_0_0_0]  (395.0,678.0) ellipse (2.0 and 2.0 );
\draw [line width=2.0pt, custcol_0_0_0 ]  (295.0,248.0) -- (394.0,248.0) ;
\draw [line width=2.0pt, custcol_0_0_0 ]  (395.0,248.0) -- (395.0,347.0) ;
\draw [line width=2.0pt, custcol_0_0_0 ]  (395.0,348.0) -- (296.0,348.0) ;
\draw [line width=2.0pt, custcol_0_0_0 ]  (295.0,348.0) -- (295.0,249.0) ;
\draw [line width=1.0pt, custcol_0_0_0 ]  (345.0,338.0) -- (343.0,348.0) -- (347.0,348.0) -- cycle;
\logisimfontD{\fontsize{14pt}{14pt}\fontseries{bx}\selectfont\node[inner sep=0, outer sep=0, custcol_0_0_0, anchor=base west] at  (375.0,303.0)  {Q};}
\logisimfontD{\fontsize{14pt}{14pt}\fontseries{bx}\selectfont\node[inner sep=0, outer sep=0, custcol_75_75_75, anchor=base west] at  (315.0,266.0)  {PR};}
\logisimfontD{\fontsize{14pt}{14pt}\fontseries{bx}\selectfont\node[inner sep=0, outer sep=0, custcol_75_75_75, anchor=base west] at  (351.0,266.0)  {CLR};}
\logisimfontD{\fontsize{14pt}{14pt}\fontseries{bx}\selectfont\node[inner sep=0, outer sep=0, custcol_0_0_0, anchor=base west] at  (303.0,303.0)  {D};}
\fill [line width=1.0pt, custcol_0_0_0]  (295.0,298.0) ellipse (2.0 and 2.0 );
\fill [line width=1.0pt, custcol_0_0_0]  (325.0,248.0) ellipse (2.0 and 2.0 );
\fill [line width=1.0pt, custcol_0_0_0]  (345.0,348.0) ellipse (2.0 and 2.0 );
\fill [line width=1.0pt, custcol_0_0_0]  (365.0,248.0) ellipse (2.0 and 2.0 );
\fill [line width=1.0pt, custcol_0_0_0]  (395.0,298.0) ellipse (2.0 and 2.0 );
\logisimfontB{\fontsize{12pt}{12pt}\fontseries{bx}\selectfont\node[inner sep=0, outer sep=0, custcol_0_0_0, anchor=base west] at  (355.0,19.0)  {RST};}
\draw [line width=2.0pt, custcol_0_0_0 ]  (351.0,5.0) -- (379.0,5.0) -- (379.0,25.0) -- (375.0,25.0) -- (365.0,29.0) -- (355.0,25.0) -- (351.0,25.0) -- cycle;
\fill [line width=2.0pt, custcol_0_0_0]  (365.0,28.0) ellipse (2.0 and 2.0 );
\logisimfontB{\fontsize{12pt}{12pt}\fontseries{bx}\selectfont\node[inner sep=0, outer sep=0, custcol_0_0_0, anchor=base west] at  (355.0,599.0)  {RST};}
\draw [line width=2.0pt, custcol_0_0_0 ]  (351.0,585.0) -- (379.0,585.0) -- (379.0,605.0) -- (375.0,605.0) -- (365.0,609.0) -- (355.0,605.0) -- (351.0,605.0) -- cycle;
\fill [line width=2.0pt, custcol_0_0_0]  (365.0,608.0) ellipse (2.0 and 2.0 );
\draw [line width=2.0pt, custcol_0_0_0 ]  (295.0,58.0) -- (394.0,58.0) ;
\draw [line width=2.0pt, custcol_0_0_0 ]  (395.0,58.0) -- (395.0,157.0) ;
\draw [line width=2.0pt, custcol_0_0_0 ]  (395.0,158.0) -- (296.0,158.0) ;
\draw [line width=2.0pt, custcol_0_0_0 ]  (295.0,158.0) -- (295.0,59.0) ;
\draw [line width=1.0pt, custcol_0_0_0 ]  (345.0,148.0) -- (343.0,158.0) -- (347.0,158.0) -- cycle;
\logisimfontD{\fontsize{14pt}{14pt}\fontseries{bx}\selectfont\node[inner sep=0, outer sep=0, custcol_0_0_0, anchor=base west] at  (375.0,113.0)  {Q};}
\logisimfontD{\fontsize{14pt}{14pt}\fontseries{bx}\selectfont\node[inner sep=0, outer sep=0, custcol_75_75_75, anchor=base west] at  (315.0,76.0)  {PR};}
\logisimfontD{\fontsize{14pt}{14pt}\fontseries{bx}\selectfont\node[inner sep=0, outer sep=0, custcol_75_75_75, anchor=base west] at  (351.0,76.0)  {CLR};}
\logisimfontD{\fontsize{14pt}{14pt}\fontseries{bx}\selectfont\node[inner sep=0, outer sep=0, custcol_0_0_0, anchor=base west] at  (303.0,113.0)  {D};}
\fill [line width=1.0pt, custcol_0_0_0]  (295.0,108.0) ellipse (2.0 and 2.0 );
\fill [line width=1.0pt, custcol_0_0_0]  (325.0,58.0) ellipse (2.0 and 2.0 );
\fill [line width=1.0pt, custcol_0_0_0]  (345.0,158.0) ellipse (2.0 and 2.0 );
\fill [line width=1.0pt, custcol_0_0_0]  (365.0,58.0) ellipse (2.0 and 2.0 );
\fill [line width=1.0pt, custcol_0_0_0]  (395.0,108.0) ellipse (2.0 and 2.0 );
\draw [line width=3.0pt, custcol_0_0_0 ]  (475.0,108.0) -- (565.0,108.0) -- (579.0,108.0) ;
\draw [line width=3.0pt, custcol_0_0_0 ]  (475.0,298.0) -- (505.0,298.0) -- (505.0,118.0) -- (565.0,118.0) -- (579.0,118.0) ;
\draw [line width=3.0pt, custcol_0_0_0 ]  (425.0,488.0) -- (475.0,488.0) -- (535.0,488.0) -- (535.0,128.0) -- (565.0,128.0) -- (579.0,128.0) ;
\draw [line width=3.0pt, custcol_0_0_0 ]  (365.0,628.0) -- (365.0,618.0) -- (425.0,618.0) -- (425.0,678.0) -- (565.0,678.0) -- (565.0,138.0) -- (579.0,138.0) ;
\draw [line width=5.0pt, custcol_0_0_0 ]  (579.0,137.0) -- (579.0,104.0) -- (584.0,99.0) ;
\logisimfontA{\fontsize{7pt}{7pt}\selectfont\node[inner sep=0, outer sep=0, custcol_0_0_0, anchor=base west] at  (572.0,105.0)  {3};}
\logisimfontA{\fontsize{7pt}{7pt}\selectfont\node[inner sep=0, outer sep=0, custcol_0_0_0, anchor=base west] at  (572.0,115.0)  {2};}
\logisimfontA{\fontsize{7pt}{7pt}\selectfont\node[inner sep=0, outer sep=0, custcol_0_0_0, anchor=base west] at  (572.0,125.0)  {1};}
\logisimfontA{\fontsize{7pt}{7pt}\selectfont\node[inner sep=0, outer sep=0, custcol_0_0_0, anchor=base west] at  (572.0,135.0)  {0};}
\fill [line width=5.0pt, custcol_0_0_0]  (585.0,98.0) ellipse (2.0 and 2.0 );
\fill [line width=5.0pt, custcol_0_0_0]  (565.0,108.0) ellipse (2.0 and 2.0 );
\fill [line width=5.0pt, custcol_0_0_0]  (565.0,118.0) ellipse (2.0 and 2.0 );
\fill [line width=5.0pt, custcol_0_0_0]  (565.0,128.0) ellipse (2.0 and 2.0 );
\fill [line width=5.0pt, custcol_0_0_0]  (565.0,138.0) ellipse (2.0 and 2.0 );
\logisimfontB{\fontsize{12pt}{12pt}\fontseries{bx}\selectfont\node[inner sep=0, outer sep=0, custcol_0_0_0, anchor=base west] at  (334.0,564.0)  {CLK};}
\draw [line width=2.0pt, custcol_0_0_0 ]  (331.0,552.0) -- (335.0,552.0) -- (345.0,548.0) -- (355.0,552.0) -- (359.0,552.0) -- (359.0,571.0) -- (331.0,571.0) -- cycle;
\fill [line width=2.0pt, custcol_0_0_0]  (345.0,548.0) ellipse (2.0 and 2.0 );
\logisimfontB{\fontsize{12pt}{12pt}\fontseries{bx}\selectfont\node[inner sep=0, outer sep=0, custcol_0_0_0, anchor=base west] at  (355.0,219.0)  {RST};}
\draw [line width=2.0pt, custcol_0_0_0 ]  (351.0,205.0) -- (379.0,205.0) -- (379.0,225.0) -- (375.0,225.0) -- (365.0,229.0) -- (355.0,225.0) -- (351.0,225.0) -- cycle;
\fill [line width=2.0pt, custcol_0_0_0]  (365.0,228.0) ellipse (2.0 and 2.0 );
\end{tikzpicture}

}
    \vspace{-35pt}
    \caption{Lógica interna do Registrador de Deslocamento}
    \label{fig:shiftRegister}
\end{figure}

\begin{table}[H]
    \centering
    \resizebox{\textwidth}{!}{
    \begin{tabular}{|c|c|c|c|c|c|c|c|c|}
        \hline
        \rowcolor{black}
        \textcolor{white}{\textbf{Estado}} & \multicolumn{8}{|c|}{\textcolor{white}{\textbf{Próximo Estado}}} \\ \cline{2-9}
        \rowcolor{black}
        \textcolor{white}{\textbf{Atual}} & \textcolor{white}{\textbf{\scriptsize 0000011}} &
        \textcolor{white}{\textbf{\scriptsize 0100011}} &
        \textcolor{white}{\textbf{\scriptsize 0010011}} &
        \textcolor{white}{\textbf{\scriptsize 1100111}} &
        \textcolor{white}{\textbf{\scriptsize 0110011}} &
        \textcolor{white}{\textbf{\scriptsize 1100011}} &
        \textcolor{white}{\textbf{\scriptsize 1101111}} &
        \textcolor{white}{\textbf{\scriptsize 0110111}} \\ \hline
         
        0000 & \cod{SL(1)} & \cod{SL(1)} & \cod{SL(1)} & \cod{SL(1)} & \cod{SL(1)} & \cod{SL(1)} & \cod{SL(1)} & \cod{SL(1)} \\ \hline

        \rowcolor{gray!20}
        0001 & \cod{SL(1)} & \cod{SL(1)} & \cod{SL(1)} & \cod{SL(1)} & \cod{SL(0)} & \cod{SR(1)} & \cod{LOAD(0101)} & \cod{LOAD(1101)} \\ \hline

        0011 & \cod{SL(0)} & \cod{SL(1)} & \cod{SR(1)} & \cod{LOAD(0101)} & X & X & X & X \\ \hline

        \rowcolor{gray!20}
        0010 & X & X & X & X & \cod{SR(1)} & X & X & X \\ \hline
        
        1000 & X & X & X & X & X & \cod{SL(0)} & X & X \\ \hline

        \rowcolor{gray!20}
        0101 & X & X & X & \cod{LOAD(0000)} & X & X & \cod{LOAD(0000)} & X \\ \hline

        0110 & \cod{SR(1)} & X & X & X & X & X & X & X \\ \hline

        \rowcolor{gray!20}
        0111 & X & \cod{LOAD(0000)} & X & X & X & X & X & X \\ \hline

        1001 & X & X & \cod{LOAD(0000)} & X & \cod{LOAD(0000)} & X & X & X \\ \hline

        \rowcolor{gray!20}
        1011 & X & \cod{LOAD(0000)} & X & X & X & X & X & X \\ \hline

        1101 & X & X & X & X & X & X & X & \cod{LOAD(0000)} \\ \hline
    \end{tabular}
    }
    \caption{Tabela de transição de estados do controle}
    \label{tab:tran_estados_abst}
\end{table}

\noindent Note que os estados foram escolhidos \textit{a priori} de forma a minimizar os `LOAD's necessários - claro que não há garantia de que essa minimização foi ótima, ela foi feita por inspeção.

Agora, basta utilizar um multiplexador de tamanho adequado para selecionar as entradas corretas para o registrador de deslocamento em cada caso, isto é, os valores de \cod{OP} (\cod{S}, na \autoref{fig:shiftRegister}), \cod{L}, \cod{R} e \cod{D} - o trem de bits que entra nos multiplexadores foi montado nessa ordem.

\subsection{Tarefa IV}
\begin{tcolorbox}[title=Enunciado, colback=blue!5!white, colframe=blue!75!black]


\end{tcolorbox}


\subsubsection{Tarefa V}
\begin{tcolorbox}[title=Enunciado, colback=blue!5!white, colframe=blue!75!black]


\end{tcolorbox}


\subsubsection{Tarefa VI}
\begin{tcolorbox}[title=Enunciado, colback=blue!5!white, colframe=blue!75!black]


\end{tcolorbox}


\subsection{Tarefa VII}
\begin{tcolorbox}[title=Enunciado, colback=blue!5!white, colframe=blue!75!black]


\end{tcolorbox}

\begin{tcolorbox}[title=Enunciado, colback=blue!5!white, colframe=blue!75!black]


\end{tcolorbox}


\begin{tcolorbox}[title=Enunciado, colback=blue!5!white, colframe=blue!75!black]


\end{tcolorbox}


\begin{tcolorbox}[title=Enunciado, colback=blue!5!white, colframe=blue!75!black]

\end{tcolorbox}



\begin{tcolorbox}[title=Enunciado, colback=blue!5!white, colframe=blue!75!black]


\end{tcolorbox}
\end{document}